% ============================================================
% Context Management in Large Language Models:
% A Survey of Human-Inspired Memory Architectures
% ============================================================
% Author: J.C. Vaught
% Date: January 2026
% ============================================================

\documentclass[11pt,letterpaper]{article}

% --- Core Packages ---
\usepackage[utf8]{inputenc}
\usepackage[T1]{fontenc}
\usepackage[margin=1in]{geometry}
\usepackage{setspace}
\usepackage{parskip}
\usepackage{xcolor}
\usepackage{titlesec}
\usepackage{hyperref}
\usepackage[numbers,sort&compress]{natbib}
\usepackage{graphicx}
\usepackage{amsmath,amssymb}
\usepackage{booktabs}
\usepackage{longtable}
\usepackage{array}
\usepackage{tabularx}
\usepackage{fancyhdr}
\usepackage{caption}
\usepackage{float}

% --- Brand Colors ---
\definecolor{garnet}{RGB}{115,0,10}
\definecolor{atlantic}{RGB}{70,106,159}
\definecolor{congaree}{RGB}{31,65,77}
\definecolor{black90}{RGB}{54,54,54}
\definecolor{black70}{RGB}{92,92,92}
\definecolor{black50}{RGB}{162,162,162}
\definecolor{rose}{RGB}{204,46,64}

% --- Section Formatting ---
\titleformat{\section}
  {\Large\bfseries\color{garnet}}
  {\thesection}{1em}{}
\titleformat{\subsection}
  {\large\bfseries\color{congaree}}
  {\thesubsection}{1em}{}
\titleformat{\subsubsection}
  {\normalsize\bfseries\color{black90}}
  {\thesubsubsection}{1em}{}
\titleformat{\paragraph}[runin]
  {\normalsize\bfseries\color{black70}}
  {}{0em}{}[.]

% --- Hyperlink Colors ---
\hypersetup{
    colorlinks=true,
    linkcolor=garnet,
    citecolor=atlantic,
    urlcolor=atlantic
}

% --- Header/Footer ---
\pagestyle{fancy}
\fancyhf{}
\fancyhead[L]{\small\color{black70}Context Management in Large Language Models}
\fancyhead[R]{\small\color{black70}J.C. Vaught}
\fancyfoot[C]{\thepage}
\renewcommand{\headrulewidth}{0.4pt}
\renewcommand{\footrulewidth}{0pt}

% --- Spacing ---
\onehalfspacing

\begin{document}

% === Title Page ===
\begin{titlepage}
\centering
\vspace*{2cm}
{\Huge\bfseries\color{garnet} Context Management in\\Large Language Models\par}
\vspace{0.75cm}
{\LARGE\color{congaree} A Survey of Human-Inspired\\Memory Architectures\par}
\vspace{2cm}
{\Large J.C. Vaught\par}
\vspace{0.5cm}
{\large January 2026\par}
\vspace{3cm}
{\large\color{black70}
A comprehensive literature review examining context window limitations, memory architectures, and human-inspired approaches to persistent memory in large language models. This survey spans attention mechanisms, positional encoding strategies, key-value cache optimization, context compression, retrieval-augmented generation, hierarchical memory systems grounded in cognitive science, hallucination phenomena arising from context mismanagement, multi-turn dialogue coherence, and evaluation methodologies for long-context capabilities.\par}
\end{titlepage}

% === Abstract ===
\begin{abstract}
\noindent
Large language models have demonstrated remarkable capabilities across diverse natural language processing tasks, yet they remain fundamentally constrained by finite context windows that limit the amount of information available during inference. This survey provides a comprehensive examination of context management strategies, organized around ten interconnected themes. We begin with the foundational attention mechanism and its quadratic computational complexity, then trace the evolution of positional encoding methods---from absolute embeddings through rotary position encodings and their scaling variants---that have progressively extended effective context lengths from thousands to millions of tokens. We analyze sparse and efficient attention architectures, including Longformer, BigBird, and state-space models such as Mamba, that achieve sub-quadratic complexity while preserving modeling quality. The survey examines key-value cache optimization techniques encompassing eviction policies, quantization schemes, and dynamic memory management systems such as PagedAttention and vAttention. Context compression methods, including prompt compression frameworks like LLMLingua and learned compression via gist tokens and in-context autoencoders, are analyzed for their trade-offs between information preservation and computational efficiency. We provide an extensive treatment of memory-augmented architectures spanning classical neural memory networks, retrieval-augmented generation and its modern variants (Self-RAG, CRAG, GraphRAG), and operating-system-inspired designs such as MemGPT. The core of this survey is dedicated to hierarchical, human-inspired memory architectures, where we draw systematic parallels between cognitive science constructs---the Atkinson-Shiffrin model, Baddeley's working memory, Tulving's memory taxonomy, Ebbinghaus's forgetting curve, and McClelland's complementary learning systems---and their computational analogues in systems such as Generative Agents, MemoryBank, the Cognitive Architectures for Language Agents (CoALA) framework, the Self-Controlled Memory framework, and the Hierarchical Memory Transformer. We examine how hallucination phenomena arise from context management failures, including the lost-in-the-middle effect and positional biases. Multi-turn dialogue presents unique challenges for memory persistence, and we survey approaches from conversation summarization to episodic memory buffers. Finally, we critically assess evaluation methodologies including Needle-in-a-Haystack, RULER, LongBench, InfiniteBench, HELMET, and domain-specific benchmarks, identifying gaps between benchmark performance and real-world long-context utilization. Throughout, we argue that human cognitive memory architectures provide essential design principles for building language model systems capable of coherent, persistent, and scalable context management.
\end{abstract}

\newpage

% ============================================================
% SECTION 1: ATTENTION FUNDAMENTALS
% ============================================================
\section{Attention Mechanisms and Their Computational Foundations}
\label{sec:attention}

The transformer architecture, introduced by Vaswani et al.\ \cite{vaswani2017attention}, fundamentally transformed natural language processing by replacing recurrent computation with a self-attention mechanism that enables direct pairwise interaction between all positions in a sequence. This architectural innovation eliminated the sequential bottleneck inherent in recurrent neural networks, enabling massive parallelization during training and establishing the foundation upon which all modern large language models are built. However, the very mechanism that grants transformers their expressive power---computing attention scores between every pair of tokens---introduces a quadratic computational complexity that constitutes the central constraint motivating the entire field of context management research surveyed in this paper.

Understanding the attention mechanism and its limitations requires examining several interconnected dimensions: the mathematical formulation that gives rise to quadratic complexity, the multi-head extension that introduces both representational richness and exploitable redundancy, the positional encoding schemes that determine sequence length generalization, and the empirically observed phenomena---the lost-in-the-middle effect and attention sinks---that reveal fundamental limitations in how transformers utilize long contexts. Each of these dimensions has spawned substantial research programs aimed at overcoming specific bottlenecks, and together they motivate the broader landscape of context management techniques surveyed in subsequent sections.

\subsection{Self-Attention and Quadratic Complexity}

The self-attention mechanism operates by projecting input tokens into three distinct representational spaces: queries, keys, and values. Given an input sequence of $n$ tokens represented as a matrix $X \in \mathbb{R}^{n \times d}$, the mechanism computes learned linear projections $Q = XW_Q$, $K = XW_K$, and $V = XW_V$, where $W_Q, W_K, W_V \in \mathbb{R}^{d \times d_k}$ are parameter matrices. The attention output is then computed as:

\begin{equation}
\text{Attention}(Q, K, V) = \text{softmax}\!\left(\frac{QK^\top}{\sqrt{d_k}}\right) V
\label{eq:attention}
\end{equation}

The matrix multiplication $QK^\top$ produces an $n \times n$ attention score matrix, requiring $O(n^2 d_k)$ floating-point operations. This quadratic dependence on sequence length $n$ means that doubling the context window quadruples the computational cost. Formal analysis by Duman-Keles et al.\ has established that the time complexity of self-attention is necessarily quadratic in the input length, with rigorous proofs showing that this constraint holds even for approximate attention mechanisms unless the Strong Exponential Time Hypothesis is false. The practical ramifications are severe: processing a 128,000-token context requires approximately 16,384 times more attention computation than processing a 1,000-token context, creating an exponential barrier to extending context windows through brute-force scaling alone.

\subsubsection{Memory Complexity and Its Decoupling from Time Complexity}

The memory complexity compounds the computational challenge. The full $n \times n$ attention matrix must be materialized during the forward pass in naive implementations, consuming $O(n^2)$ memory. For a sequence of 100,000 tokens with 16-bit precision, this attention matrix alone requires approximately 18.6 gigabytes per attention layer, rendering naive full-attention inference infeasible on current hardware for the context lengths demanded by practical applications such as book-length document analysis or multi-session dialogue.

Importantly, subsequent research demonstrated that $O(n^2)$ memory is not strictly necessary even when time complexity remains quadratic. Rabe and Staats showed that self-attention does not require $O(n^2)$ memory by computing attention in blocks that avoid materializing the full attention matrix in global memory. The key insight is that the softmax normalization can be decomposed into partial computations that are combined without storing intermediate full-size matrices. This decoupling of memory from time complexity laid the theoretical groundwork for FlashAttention \cite{dao2022flashattention}, discussed in Section~\ref{sec:sparse}, which exploits this principle with hardware-aware tiling strategies.

\subsubsection{Computational Cost at Scale}

To appreciate the severity of the quadratic bottleneck in practical terms, consider the progression of context windows in modern language models. GPT-2 in 2019 operated with a context window of 1,024 tokens; GPT-3 in 2020 expanded to 2,048 tokens; GPT-3.5 in 2022 reached 4,096 tokens; and GPT-4 Turbo in 2023 achieved 128,000 tokens. This represents a roughly 128$\times$ increase in context length, but due to the quadratic scaling of attention, the computational cost of the attention mechanism alone increased by a factor of approximately 16,384$\times$. Without efficiency innovations, this scaling trajectory would be computationally unsustainable. The Gemini 1.5 model family, released in 2024 with context windows of 1,000,000 tokens, would require approximately $10^{12}$ attention score computations per layer under naive full attention---a figure that underscores the absolute necessity of the efficiency techniques surveyed in Sections~\ref{sec:sparse} and \ref{sec:kvcache}.

The inference-time cost structure differs importantly from training-time costs. During autoregressive generation, the model processes one token at a time, with each new token requiring attention computation against all previous tokens in the sequence. This creates an $O(n)$ cost per generated token that accumulates to $O(n^2)$ over the full sequence, but the per-step memory-bandwidth cost of loading the KV cache from GPU memory becomes the dominant bottleneck in practice. A model with 8 billion parameters may require loading 280 gigabytes of KV cache during generation---30$\times$ larger than the model parameters themselves---with inference efficiency dropping from approximately 70\% during training to roughly 10\% during generation due to memory-bandwidth limitations. This memory-bandwidth bottleneck motivates the KV cache optimization techniques discussed in Section~\ref{sec:kvcache}.

\subsection{Multi-Head Attention and Head Specialization}

The original transformer architecture extends scalar attention to multi-head attention (MHA), where $h$ independent attention heads operate in parallel, each with its own projection matrices:

\begin{equation}
\text{MultiHead}(Q, K, V) = \text{Concat}(\text{head}_1, \ldots, \text{head}_h) W_O
\end{equation}

where $\text{head}_i = \text{Attention}(QW_Q^i, KW_K^i, VW_V^i)$ and $W_O \in \mathbb{R}^{hd_v \times d}$ is an output projection. This formulation enables the model to attend to information from different representation subspaces at different positions simultaneously \cite{vaswani2017attention}. The total computational cost remains $O(n^2 d)$ because the head dimension $d_k = d/h$ decreases proportionally with the number of heads, maintaining constant total computation while distributing it across specialized sub-computations.

\subsubsection{Functional Specialization of Attention Heads}

Empirical analysis has revealed that attention heads exhibit meaningful functional specialization that emerges naturally during pre-training. Studies categorizing heads across multiple model architectures have identified distinct functional types. Positional heads focus on local position patterns, attending predominantly to adjacent tokens and encoding sequential structure. Syntactic heads capture grammatical relationships such as subject-verb dependencies, relative clause attachments, and prepositional phrase bindings. Rare or specialized heads activate selectively for specific linguistic phenomena such as coreference resolution, negation scope, or quotation attribution.

This specialization varies systematically across layers. Early layers tend toward more balanced, preprocessing-oriented head usage, performing operations analogous to tokenization refinement and basic syntactic parsing. Middle layers exhibit mixed utilization patterns, with some heads specializing in semantic relationships and others maintaining syntactic tracking. Late layers are often dominated by one or two heads responsible for output generation, implementing a kind of attention bottleneck that concentrates the model's final prediction on the most relevant information. Further analysis has demonstrated that in the context of in-context learning, first-layer heads preprocess examples while subsequent layers employ a single dominant head for iterative optimization, effectively implementing a preprocess-then-optimize algorithm within the forward pass.

The layer-wise variation in head utilization has important implications for cache management strategies discussed in Section~\ref{sec:kvcache}. Because early layers tend to distribute attention more broadly while later layers concentrate attention on fewer tokens, the optimal cache eviction strategy may differ substantially across layers---an insight exploited by FastGen \cite{ge2024model} and PyramidKV \cite{cai2024pyramidkv}, which apply different cache budgets and eviction policies per layer.

\subsubsection{Head Redundancy and Efficient Attention Variants}

A significant finding with direct implications for context management efficiency is the extensive redundancy among attention heads. The Mixture-of-Head (MoH) attention framework demonstrated through experiments on Vision Transformers, Diffusion Transformers, and large language models that MoH outperforms standard multi-head attention while using only 50--90\% of available heads. Research on attention head pruning has shown that the vast majority of heads---especially encoder self-attention heads---can be removed with minimal impact on task performance, suggesting that standard multi-head attention over-provisions representational capacity.

This redundancy has been exploited commercially through two important architectural innovations. Grouped Query Attention (GQA) \cite{ainslie2023gqa} shares key-value projections across groups of query heads, reducing the key-value cache size by a factor proportional to the group size. GQA achieves 10--100$\times$ smaller KV caches and up to 12$\times$ faster decoder inference with minimal quality degradation, and has been adopted in production models including Llama~2, Llama~3, and Gemma. For the Llama~3 8B model, GQA reduces KV cache memory from 280 gigabytes to 69 gigabytes. Multi-Query Attention (MQA) \cite{shazeer2019fast} represents the extreme case where a single key-value head serves all query heads, achieving maximum cache compression at a modest quality cost.

The progression from MHA to GQA to MQA illustrates a general principle in context management: exploiting redundancy in intermediate representations enables substantial efficiency gains with tolerable quality degradation. This principle recurs throughout the survey, appearing in cache compression (Section~\ref{sec:kvcache}), context compression (Section~\ref{sec:compression}), and memory hierarchies (Section~\ref{sec:hierarchical}).

\subsection{Positional Encoding and Sequence Length Generalization}

Because the self-attention operation in Equation~\ref{eq:attention} is permutation-equivariant---it produces the same output regardless of token ordering absent positional information---transformers require explicit positional encodings to capture sequence structure. The choice of positional encoding scheme profoundly affects the model's ability to generalize to sequence lengths beyond those encountered during training, making positional encoding one of the most actively researched topics in the context extension literature.

\subsubsection{Absolute Positional Embeddings}

The original transformer employed fixed sinusoidal positional encodings defined by sine and cosine functions of different frequencies:

\begin{equation}
PE_{(pos, 2i)} = \sin\!\left(\frac{pos}{10000^{2i/d}}\right), \qquad PE_{(pos, 2i+1)} = \cos\!\left(\frac{pos}{10000^{2i/d}}\right)
\end{equation}

where $pos$ is the position index and $i$ is the dimension index. These deterministic encodings require no learned parameters and theoretically enable the model to learn relative position information through linear combinations of the encoding vectors. The choice of sinusoidal functions provides desirable properties: the encoding for any fixed offset $k$ can be expressed as a linear function of the encoding at position $pos$, enabling the model to attend to relative positions through learned linear projections.

Subsequent models adopted learned absolute positional embeddings, where each position index receives a trainable embedding vector added to the token embedding. BERT, GPT-2, and GPT-3 all employed this approach. While learned embeddings achieve better in-distribution performance given sufficient training data, they fundamentally limit length generalization: the model encounters novel position indices at inference time when processing sequences longer than those seen during training, leading to unpredictable behavior. Sinusoidal encodings offer better extrapolation in principle, but in practice neither approach enables reliable generalization beyond the training length. This fundamental limitation motivated the development of relative and rotary encoding schemes that encode positional relationships rather than absolute positions.

\subsubsection{Relative Positional Encodings}

Shaw et al.\ \cite{shaw2018self} introduced relative positional representations that encode the distance between tokens rather than their absolute positions, achieving a 1.3 BLEU improvement on WMT 2014 English-to-German translation. By encoding pairwise distances between sequence elements, relative encodings enable the model to generalize to sequences of arbitrary length in principle, since the set of possible pairwise distances does not change with sequence length. This conceptual advance was critical because it decoupled positional information from absolute position indices, removing the hard ceiling on sequence length imposed by absolute embeddings.

Transformer-XL \cite{dai2019transformerxl} further developed the relative positional approach with two key innovations. First, it introduced learnable relative position biases that replace absolute position embeddings in the attention computation, avoiding what the authors termed ``temporal confusion'' when segment-level recurrence reuses position indices across segments. Second, it introduced segment-level recurrence, where hidden states from previous segments are cached and concatenated with the current segment's representations during attention computation. Together, these innovations extend effective context length by 80--450\% compared to vanilla transformers, with Transformer-XL modeling 80--133\% longer dependencies than RNNs and 291--447\% longer than standard transformers on language modeling benchmarks.

ALiBi \cite{press2022train} dramatically simplified relative encoding by adding a linear bias proportional to the distance between tokens directly to attention scores, with different bias slopes for different attention heads. Rather than encoding position in the embedding space, ALiBi modifies the attention computation itself: $\text{softmax}(q_i^\top k_j - m \cdot |i - j|)$, where $m$ is a head-specific slope parameter. This approach requires no additional parameters, provides native length extrapolation capability (a 1.3B parameter model trained on 1,024 tokens maintains low perplexity at 16,000 tokens), and offers 11\% faster training and 11\% less memory consumption compared to sinusoidal encodings. ALiBi has been adopted in production models including MPT and BLOOM, and its success demonstrated that position encoding need not be a separate embedding but can instead be integrated directly into the attention scoring function.

\subsubsection{Rotary Position Embedding}

Rotary Position Embedding (RoPE) \cite{su2024roformer} encodes positional information by rotating query and key vectors in a high-dimensional space, where the angle of rotation is proportional to the token's position. Mathematically, RoPE applies a rotation matrix $R_\theta(m)$ to the query and key vectors at position $m$:

\begin{equation}
q_m = R_\theta(m) W_Q x_m, \qquad k_n = R_\theta(n) W_K x_n
\end{equation}

The inner product $q_m^\top k_n$ then depends only on the relative position $m - n$ through the rotation difference $R_\theta(m-n)$. The rotation is applied pairwise to embedding dimensions, with each pair rotated by an angle $m \cdot \theta_i$ where $\theta_i = 10000^{-2i/d}$ defines the frequency for dimension pair $i$. This creates a natural frequency decomposition: low-indexed dimension pairs rotate at high frequencies (encoding fine-grained local position), while high-indexed pairs rotate at low frequencies (encoding coarse global position).

RoPE provides several desirable properties that have made it the dominant positional encoding in modern large language models. It encodes relative position without additional parameters beyond the base frequency. It is compatible with linear self-attention approximations. Its frequency decomposition enables principled scaling strategies (discussed in Section~\ref{sec:window}), where different frequency bands can be scaled independently based on their contribution to local versus global position encoding. RoPE has been adopted by the Llama, Qwen, DeepSeek, Mistral, and Gemma model families, making its scaling properties central to the entire context extension research program.

However, RoPE suffers from poor extrapolation to positions beyond the training range. When position indices exceed the training maximum, the rotation angles become unfamiliar and attention scores can become numerically unstable. This limitation motivated an extensive body of work on RoPE scaling methods---Position Interpolation, NTK-aware scaling, YaRN, and LongRoPE---discussed in Section~\ref{sec:window}.

\subsubsection{Comparative Analysis of Positional Encoding Approaches}

The evolution from sinusoidal to learned to relative to rotary position encodings reflects a progressive understanding of what properties are essential for effective sequence modeling. Sinusoidal encodings are deterministic and theoretically capable of encoding relative positions, but fail to generalize reliably in practice. Learned embeddings achieve superior in-distribution performance but impose a hard length ceiling. Relative encodings (Shaw, Transformer-XL) remove the length ceiling but introduce computational overhead. ALiBi provides an elegant, parameter-free solution with native extrapolation but limited fine-grained position control. RoPE achieves the best balance of properties---relative encoding, no extra parameters, natural frequency decomposition---but requires explicit scaling for extrapolation.

Recent distributional analysis has shown that rotary angle distributions are critical for context extension success, with existing scaling methods introducing out-of-distribution angles that cause suboptimal performance. A distributional minimization approach that selectively applies interpolation or extrapolation per dimension based on disturbance scores achieves 4.33\% improvement over Position Interpolation when extending Llama~2 to 16K tokens. This analysis suggests that optimal position encoding design requires understanding not just the functional form but also the statistical properties of the rotation angles encountered during training and inference.

\subsection{The Historical Trajectory of Context Window Scaling}

The rapid growth of context windows provides essential background for understanding the urgency and importance of context management research. Context windows have grown approximately 1,000-fold in five years, from GPT-2's 1,024 tokens in 2019 to Gemini 1.5's 1,000,000 tokens in 2024 and Gemini 2.5's 2,000,000 tokens announced in 2025. This exponential growth trajectory has been enabled by the confluence of positional encoding innovations (RoPE scaling), efficiency improvements (FlashAttention, sparse attention), hardware advances (larger GPU memory, multi-device inference), and training methodology refinements (progressive context extension, long-context alignment).

Key inflection points include GPT-4 Turbo's 128K context in 2023, which demonstrated that six-figure token contexts were commercially viable, and the simultaneous achievement of million-token contexts by multiple model families in 2024. The open-source ecosystem tracked this progression: Llama~2 launched with 4,096 tokens but spawned extensive community work on context extension, while Llama~3.1 natively supported 128,000 tokens, matching proprietary models.

However, this growth in nominal context length has significantly outpaced growth in effective context utilization. Empirical evidence consistently shows that models claiming 128K or larger context windows exhibit significant performance degradation on tasks requiring genuine long-context understanding. The RULER benchmark revealed that models claiming 32K context maintain acceptable performance on only about 50\% of challenging tasks. InfiniteBench showed that accuracy on genuine 100K-token tasks remains between 20--50\%. BABILong found that effective context utilization is often only 10--20\% of the nominal window. This gap between nominal and effective context length is the central challenge motivating the evaluation methodologies discussed in Section~\ref{sec:evaluation} and the memory architectures proposed in Section~\ref{sec:hierarchical}.

\subsection{The Lost-in-the-Middle Phenomenon}

Liu et al.\ \cite{liu2024lost} documented a striking failure mode of language models processing long contexts: when relevant information is placed in the middle of a long input, model performance degrades substantially compared to when the same information appears at the beginning or end. Across models including GPT-3.5 with 4K and 16K context windows, models exhibited a characteristic U-shaped performance curve where accuracy was highest for information at the first and last positions and lowest for information near the center of the context. The performance degradation from boundary to middle positions reached as high as 63\%, demonstrating that the effective context for information retrieval is substantially smaller than the nominal context window.

\subsubsection{Connection to Human Serial Position Effects}

This phenomenon has been linked to the serial position effect long studied in human cognitive psychology \cite{murdock1962serial, glanzer1966two}. Humans similarly exhibit primacy effects (better recall for items presented first) and recency effects (better recall for items presented last), producing the classic U-shaped serial position curve first characterized by Murdock in 1962. However, the mechanisms differ fundamentally. In humans, primacy effects arise from deeper encoding through rehearsal---early items receive more rehearsal cycles in working memory, strengthening their encoding into long-term memory. Recency effects reflect the continued availability of recent items in short-term memory \cite{atkinson1968human}, which has not yet decayed. In transformers, the U-shaped curve appears to arise from attention distribution patterns, particularly the attention sink phenomenon discussed below, and from the structure of causal attention masks that give initial tokens privileged access to subsequent computation because they have been included in more attention computations than later tokens.

\subsubsection{Causes and Contributing Factors}

The lost-in-the-middle effect is not solely attributable to positional encoding limitations. Research has identified multiple contributing factors. First, position frequency distributions during training create an undertraining effect for long-distance positions: in standard pre-training, initial positions appear in every training example while later positions appear only in longer examples, creating a frequency bias toward early positions. Second, softmax crowding at long distances causes attention scores to become numerically similar, reducing the model's ability to discriminate between relevant and irrelevant tokens in the middle of long contexts. Third, the causal attention mask means that initial tokens accumulate the most attention computation across the sequence, giving them a computational advantage in shaping downstream representations.

\subsubsection{Implications for Retrieval-Augmented Generation}

The lost-in-the-middle effect has profound implications for retrieval-augmented generation systems, where the ordering of retrieved documents within the prompt directly affects answer quality. Even models with 16K context windows showed significant performance degradation when key information was positioned in the middle 50\% of the context. This has motivated practical mitigation strategies: placing the most relevant retrieved documents at the beginning and end of the context (exploiting boundary effects), reordering documents based on query relevance and position sensitivity, and designing retrieval systems that account for positional biases in the downstream language model. LongLLMLingua \cite{jiang2024longllmlingua}, discussed in Section~\ref{sec:compression}, explicitly addresses this through document reordering during prompt compression.

The phenomenon also extends beyond text to multimodal settings. Research presented at NAACL 2025 demonstrated that similar patterns of degraded performance for middle-positioned information appear in vision-language models processing long multimodal contexts, suggesting that the lost-in-the-middle effect is a fundamental property of the transformer attention mechanism rather than a text-specific artifact.

\subsection{The Attention Sink Phenomenon}

Xiao et al.\ \cite{xiao2024efficient} identified and characterized the attention sink phenomenon, wherein language models allocate disproportionately large attention scores to initial tokens---particularly the beginning-of-sequence token---even when these tokens are semantically irrelevant to the current prediction. Empirical measurement reveals that 30--50\% of total attention weight concentrates on the first few tokens across diverse models and inputs.

\subsubsection{Mechanistic Explanation}

The mechanism underlying attention sinks is rooted in the softmax normalization constraint: because softmax attention weights must sum to one, the model must allocate its attention budget somewhere. When no position is strongly semantically relevant, the model defaults to early positions as recipients of this excess attention, effectively using them as ``dump'' locations for the attention that softmax requires to be distributed. This pattern emerges during pre-training as an adaptation to the normalization constraint and compounds across layers: attention concentration in Layer~1 propagates to Layer~2 and beyond, creating a cascade effect that effectively starves middle-context tokens of attention resources.

The attention sink phenomenon is universal across transformer architectures. Research published at ICLR 2025 provided extensive empirical investigation of the precise conditions under which attention sinks emerge, analyzing which specific tokens and layers exhibit the strongest sink behavior. The study found that the phenomenon is particularly pronounced in layers with broader attention patterns (those responsible for global information aggregation rather than local processing) and that certain tokens---particularly the beginning-of-sequence token and the first semantically meaningful token---serve as sinks more consistently than others. The cascade effect means that even a modest concentration of attention on initial tokens in early layers produces severe attention starvation for middle-position tokens in later layers.

\subsubsection{StreamingLLM}

The StreamingLLM framework, proposed alongside the attention sink analysis \cite{xiao2024efficient}, demonstrates that attention sinks can be harnessed constructively. By retaining the initial sink tokens in the KV cache alongside a sliding window of recent tokens, StreamingLLM enables models trained on finite contexts to process sequences of effectively infinite length without fine-tuning. The framework achieves processing of 4~million or more tokens with up to 22.2$\times$ speedup over sliding window baselines across Llama-2, Falcon, Pythia, and MPT architectures. The key insight is that attention sinks serve as anchor points for the attention distribution, and removing these initial tokens from the KV cache causes catastrophic quality degradation, while retaining just 4--8 initial tokens alongside recent tokens preserves generation quality.

The success of StreamingLLM highlights a broader principle in context management: understanding the internal dynamics of the attention mechanism---not just its mathematical formulation---enables qualitatively new system designs. By recognizing that certain tokens serve a structural rather than semantic role, StreamingLLM achieves efficient streaming inference through a simple architectural insight rather than complex algorithmic innovation.

\subsubsection{Implications for Context Management}

The attention sink phenomenon has several important implications for context management. First, it explains a significant component of the lost-in-the-middle effect: middle tokens receive minimal attention precisely because initial tokens absorb a disproportionate share of the attention budget. Second, it constrains KV cache eviction strategies: early tokens cannot be evicted from cache without degrading model quality, even if they carry no semantic information relevant to the current query. This constraint must be respected by eviction-based cache compression methods such as H$_2$O \cite{zhang2024h2o} and Scissorhands \cite{liu2023scissorhands}, discussed in Section~\ref{sec:kvcache}. Third, it suggests that architectural modifications---such as training with explicit placeholder sink tokens, adopting alternative attention mechanisms that do not require softmax normalization, or implementing linear attention variants---may be necessary to fully utilize long contexts. The exploration of such alternatives is discussed in Section~\ref{sec:sparse}.


% ============================================================
% SECTION 2: WINDOW EXTENSION
% ============================================================
\section{Context Window Extension Techniques}
\label{sec:window}

The practical utility of large language models is directly constrained by their effective context window---the maximum number of tokens that can be processed while maintaining generation quality. Early transformer models operated with context windows of 512--2,048 tokens, sufficient for short documents but inadequate for applications requiring processing of books, codebases, extended conversations, or large document collections. The past three years have witnessed rapid progress in extending context windows from thousands to millions of tokens through innovations in positional encoding scaling, training methodology, and distributed computation. This section surveys the principal techniques that have driven this expansion, organized around four main themes: RoPE scaling methods, continual pre-training and alignment strategies, distributed attention mechanisms, and emerging hybrid approaches.

\subsection{RoPE Scaling Methods}

Because RoPE \cite{su2024roformer} has become the dominant positional encoding in modern large language models, methods for extending context beyond the training length have focused on manipulating RoPE's frequency structure. The central challenge is that RoPE encodes position through rotations at specific frequencies, and extrapolating to positions beyond the training range causes attention scores to blow up as the model encounters rotation angles it has never observed. The progression from Position Interpolation through NTK-aware scaling to YaRN and LongRoPE represents increasing sophistication in addressing this challenge, with each method building on insights from its predecessors to achieve greater extension ratios with less computational cost.

\subsubsection{Position Interpolation}

Chen et al.\ \cite{chen2023extending} proposed Position Interpolation (PI), which linearly down-scales input position indices to fit within the original context window rather than extrapolating beyond it. For a model trained on context length $L$ applied to a sequence of length $L' > L$, PI maps position $m$ to $m \cdot L / L'$, compressing all positions into the range $[0, L]$ where the model has been trained. This conceptually simple approach achieves 32$\times$ context extension on Llama models with minimal fine-tuning. Theoretical analysis demonstrates that the upper bound of interpolation error is approximately 600$\times$ smaller than that of direct extrapolation, explaining the superior stability of interpolation-based approaches.

However, Position Interpolation applies uniform scaling across all frequency dimensions, which introduces a fundamental limitation. High-frequency components of RoPE encode fine-grained local position information that should not be aggressively compressed, as doing so distorts the model's ability to distinguish adjacent tokens. Low-frequency components encode coarse global position information that can tolerate more compression, since they primarily distinguish distant positions. Uniform compression over-smooths local position information while under-compressing global position information, leading to suboptimal performance at both local and global scales. This observation---that different frequency bands should be scaled differently---motivated the development of frequency-aware scaling methods.

\subsubsection{NTK-Aware Scaling}

NTK-aware scaling addresses Position Interpolation's uniform scaling limitation by modifying the RoPE base frequency using the formula $\text{base}_\text{new} = \text{base} \times \alpha^{d/(d-2)}$, where $\alpha$ is the extension factor and $d$ is the embedding dimension. This formulation draws on insights from Neural Tangent Kernel theory to spread interpolation pressure across dimensions non-uniformly: high-frequency components encoding local patterns are scaled less aggressively, while low-frequency components encoding global position absorb more of the compression. The method enables 8K or more context extension without any fine-tuning through a modification requiring only three lines of code, making it remarkably practical for deployment.

The theoretical motivation from Neural Tangent Kernel theory is that the effective learning dynamics of neural networks depend on the frequency spectrum of the input representation. By adjusting the base frequency, NTK-aware scaling modifies the frequency spectrum to maintain the relative importance of different frequency bands even at extended lengths. The practical simplicity of this approach---a single-line code change to the RoPE base frequency---has made it widely adopted in open-source implementations and has inspired a family of dynamic variants discussed below.

\subsubsection{YaRN}

YaRN (Yet Another RoPE extensioN) \cite{peng2024yarn} combines NTK-by-parts interpolation with attention temperature scaling, achieving state-of-the-art context extension with 10$\times$ fewer tokens and 2.5$\times$ fewer training steps than Position Interpolation. The NTK-by-parts approach divides RoPE dimensions into three categories based on their wavelength relative to the training context length: high-frequency dimensions whose wavelength is much shorter than the training length are not interpolated at all (preserving local position encoding fidelity), low-frequency dimensions whose wavelength exceeds the training length are fully interpolated (compressing global position as needed), and a transition band where interpolation is applied with a smooth ramp function that gradually increases the interpolation factor.

Additionally, YaRN applies a learned temperature scaling factor to the attention logits that corrects for the distributional shift introduced by interpolation. Because interpolation compresses position information, the variance of attention logits decreases, causing softmax outputs to become more uniform and reducing the model's ability to attend selectively. The temperature correction restores the sharpness of attention distributions, maintaining the model's discriminative capacity.

YaRN enables Llama~2 to extrapolate to 128K tokens using only 0.1\% of the original pre-training data for fine-tuning, representing a dramatic improvement in data efficiency. The method has been widely adopted in production systems including Qwen, DeepSeek, Llama~3, and various open-source implementations, establishing it as the de facto standard for RoPE-based context extension. The success of YaRN demonstrates that principled, frequency-aware manipulation of positional encoding can achieve the same extension as brute-force retraining at a fraction of the computational cost.

\subsubsection{LongRoPE}

LongRoPE \cite{ding2024longrope} pushes context extension to 2,048K tokens through three complementary innovations. First, it identifies and exploits non-uniformities in the optimal interpolation factors across RoPE dimensions, using an evolutionary search algorithm to find dimension-specific scaling factors rather than relying on the analytic formulas of NTK-aware scaling or YaRN. This search-based approach discovers scaling configurations that no closed-form method captures, achieving better performance at extreme extension ratios. Second, it employs a progressive extension strategy that first extends to 256K tokens and then uses the resulting model as initialization for further extension to 2,048K, avoiding the instability of extreme single-step extension. This progressive approach mirrors curriculum learning strategies and prevents the catastrophic loss of short-context performance that can occur with aggressive single-step extension. Third, a readjustment phase recovers short-context performance that may degrade during long-context fine-tuning by applying a modified set of scaling factors optimized for shorter contexts. LongRoPE achieves 2-million-token contexts with only 1,000 fine-tuning steps, published at ICML~2024.

\subsubsection{Dynamic NTK Scaling}

Dynamic NTK scaling addresses the practical limitation of fixed scaling factors by adapting the scaling factor at inference time based on the actual sequence length being processed. When the input sequence is shorter than the training context, minimal or no scaling is applied; as the sequence approaches and exceeds the training length, scaling increases proportionally. This approach eliminates the need to choose a fixed extension factor at deployment time and allows a single model deployment to handle both short and long sequences optimally. The adaptive nature of dynamic NTK scaling mirrors the concept of adaptive memory management that recurs throughout this survey: rather than committing to fixed resource allocation, the system adjusts its strategy based on the demands of the current input.

\subsubsection{Emerging RoPE Scaling Approaches}

Several recent methods extend beyond the NTK-YaRN paradigm. CLEX (Continuous Length Extrapolation) models the dynamics of length scaling factors through ordinary differential equations, enabling smooth context extension to 4--8$\times$ the training length without performance degradation. A model trained on 4K tokens achieves competitive 32K performance on LongBench through CLEX, with minimal training and inference latency impact. E$^2$-LLM (Efficient and Extreme Length Extension) achieves extreme extension through single efficient training on short sequences using a dual RoPE augmentation strategy that scales and adjusts position indices across training samples, dramatically reducing the need for long-context training data. Segmented base adjustment proposes segment-wise adjustment of RoPE base frequencies rather than global scaling, applying optimized adjustments per input segment for more flexible context extension. Decomposed positional vector analysis decomposes positional vectors into frequency-dependent components, enabling per-frequency decisions about whether to interpolate or extrapolate, advancing understanding of how different frequency bands should be handled independently.

\subsection{Continual Pre-Training and Long-Context Alignment}

Extending the context window through positional encoding modifications is necessary but not sufficient for effective long-context utilization. Research has demonstrated that models require both the capacity to process longer sequences (provided by RoPE scaling) and the training experience to effectively utilize longer contexts (provided by continual pre-training and alignment). This subsection examines training strategies that complement positional encoding modifications.

\subsubsection{Continual Pre-Training for Context Extension}

LongAlign \cite{longalign2024} provides a comprehensive recipe for long-context alignment extending to 64K tokens through three components: expanding the RoPE base frequency by a factor of 200$\times$ (from 10,000 to 2,000,000), conducting continual training on sequences up to 64K tokens for 10 billion tokens, and developing loss-weighting strategies for packed training with varied sequence lengths. LongAlign outperforms existing recipes by up to 30\% on long-context tasks while maintaining short-context proficiency through carefully engineered data construction and training procedures. This demonstrates that without explicit alignment on long-context tasks, models with extended context windows may fail to attend to relevant information spread throughout the context, underscoring a distinction between nominal context length and effective context length.

LongSkywork provides an end-to-end training recipe for efficient context extension through dynamic batching with sorted sequences, efficient packed training with loss weighting, context window scheduling, and instruction tuning. Code Llama demonstrates domain-specific context extension through continual pretraining on code sequences, leveraging code's structural properties (syntactic nesting, cross-file references) for extended context utilization, achieving 8K--100K context windows.

\subsubsection{Data Engineering for Long-Context Training}

A critical insight from the context extension literature is that data quality and diversity are as important as training methodology. Research on data engineering for scaling to 128K contexts demonstrates that strategic data selection outperforms naive long-sequence collection approaches. The optimal training mixture balances long-context documents (providing length experience) with short-context documents (preventing degradation on short tasks), with curriculum-based sequence length scheduling---progressively increasing context length during training---proving more efficient than static long-context training.

Counter-intuitively, research has shown that supervised fine-tuning at extended context lengths can improve short-context performance, demonstrating a bidirectional relationship between short and long-context capabilities during training. This finding suggests that long-context training does not simply extend capability but may fundamentally improve the model's ability to organize and utilize information at all scales.

\subsubsection{Sample-Efficient Context Extension}

An emerging research direction demonstrates that context window extension can be achieved with remarkably few training samples. One line of work shows that context extension is possible with as few as 100 carefully selected training samples, through principled selection of training data and optimization strategies. SkyLadder introduces context window scheduling during pretraining, progressively increasing sequence length across training phases, reducing total compute while improving final model quality. Partial context training demonstrates that efficient long-context training can be achieved through partial context masking during training, reducing memory requirements while maintaining full context capability at inference time, with counterintuitive benefits from incomplete context exposure during training.

\subsection{Distributed Attention for Extended Context}

\subsubsection{Ring Attention}

Ring Attention \cite{liu2023ring} tackles context extension at the systems level by distributing attention computation across multiple devices. The key insight is that the blockwise computation of attention can be mapped to a ring communication topology among devices. Each device holds a block of the key-value cache and circulates it around the ring, enabling each device to compute its portion of the attention output without requiring the full context to reside on any single device. Communication overhead is overlapped with computation through careful scheduling: while device $i$ is computing attention against its current block of keys, it simultaneously sends that block to device $i+1$ and receives the next block from device $i-1$.

Ring Attention enables near-infinite context lengths limited primarily by total available memory across devices rather than by individual device memory. The framework has been used to demonstrate million-token context processing on clusters of consumer-grade GPUs, making very long context accessible outside large-scale infrastructure. Ring Attention represents a qualitatively different approach to context extension---rather than making attention more efficient within a device, it distributes exact full attention across devices---and has influenced production-scale deployment strategies.

\subsection{Emerging Context Extension Paradigms}

Several emerging approaches signal future directions in context extension beyond the RoPE-scaling paradigm.

\subsubsection{Attention-Level Innovations}

Core Context Aware Attention (CCA-Attention) proposes globality-pooling attention where tokens communicate via a reduced-size core token set, reducing computational complexity while maintaining semantic understanding through pooled global representations. Adaptive Grouped Attention proposes grouping tokens for attention computation with group sizes that adapt based on position and importance, reducing context scaling complexity with minimal training overhead. MixAttention from Databricks mixes local and global attention patterns during inference, combining the efficiency of sparse local patterns with strategic global connections for practical long-context processing.

\subsubsection{Probabilistic and Theoretical Frameworks}

A Bayesian attention mechanism has been proposed that treats positional uncertainty explicitly, enabling principled length extrapolation through Bayesian inference rather than empirical scaling heuristics. This provides a theoretical foundation for understanding when and why certain extrapolation strategies succeed, moving the field from engineering solutions toward principled design. Meta-analysis has argued that scaling context length requires fundamental rethinking of attention mechanisms rather than positional encoding tweaks alone, advocating for integrated approaches combining position scaling, attention approximation, and memory strategies for optimal long-context performance.

\subsubsection{Training-Time versus Inference-Time Approaches}

A growing body of work demonstrates that training-time decisions about attention patterns fundamentally affect inference-time context handling. Sliding window attention training proposes training LLMs with sliding window attention masks from scratch, demonstrating that the choice of attention pattern during training determines the effective context at inference time. TokenSelect proposes selective token retention during inference to reduce KV cache requirements, using learned scoring to identify important tokens for retention. These approaches suggest that the distinction between ``model architecture'' and ``inference strategy'' is increasingly blurred, with the most effective systems co-designing training procedures and inference algorithms.


% ============================================================
% SECTION 3: SPARSE AND EFFICIENT ATTENTION
% ============================================================
\section{Sparse and Efficient Attention Architectures}
\label{sec:sparse}

The quadratic complexity of full self-attention has motivated a rich body of work on approximate and efficient attention mechanisms. These approaches fall into several broad categories: sparse attention patterns that restrict which token pairs interact, linear attention approximations that avoid computing the full attention matrix, hardware-aware implementations that optimize memory access patterns, and alternative architectures that replace attention entirely with sub-quadratic sequence mixing operations.

\subsection{Sparse Attention Patterns}

Sparse attention methods impose structure on the attention matrix by restricting each token to attend to only a subset of positions, reducing the number of computed attention scores from $O(n^2)$ to $O(n \cdot k)$ where $k \ll n$ is the effective attention span per token.

\subsubsection{Longformer}

The Longformer \cite{beltagy2020longformer} introduced a combination of sliding window attention and global attention. Each token attends to a fixed-size local window of neighboring tokens (typically 256--512 tokens on each side), capturing local syntactic and semantic patterns with $O(n \cdot w)$ complexity. Additionally, a small number of designated tokens receive global attention, attending to and being attended by all other tokens. This global-local hybrid enables the model to combine fine-grained local processing with coarse global information aggregation, and established the principle that full pairwise attention is unnecessary for many tasks.

\subsubsection{BigBird}

BigBird \cite{zaheer2020bigbird} extended the sparse attention framework with a theoretical foundation, proving that the combination of random attention, sliding window attention, and global attention is a universal approximator of full attention. The random attention component ensures that the sparse attention graph remains well-connected, providing the theoretical guarantee that information can flow between any two positions within a bounded number of attention layers. BigBird achieves $O(n)$ complexity while maintaining the expressiveness of full attention, as demonstrated on question answering, summarization, and genomics tasks.

\subsubsection{Sparse Transformers}

OpenAI's Sparse Transformer proposed a two-dimensional factorization of the attention matrix using strided and fixed attention patterns, achieving $O(n\sqrt{n})$ complexity. The strided pattern has each token attend to tokens at regular intervals (every $\sqrt{n}$ positions), while the fixed pattern designates specific positions as summary tokens that attend to and are attended by all positions in their block. The combination ensures that information can flow between any two positions within two attention layers. The approach effectively handles sequences of 60,000 or more tokens across modalities including text, images, and audio, demonstrating that sparse attention is viable not only for text but for general sequence modeling.

\subsection{Linear Attention Approximations}

\subsubsection{Performer}

The Performer \cite{choromanski2021rethinking} approximates the softmax attention kernel using random feature maps, decomposing the attention computation into a form that avoids explicit computation of the $n \times n$ attention matrix. The Performer replaces the exponential kernel with a positive random feature approximation $\phi(q)^\top \phi(k)$, enabling the attention output to be computed as $(\phi(Q) (\phi(K)^\top V))$ rather than $((\phi(Q) \phi(K)^\top) V)$, reducing complexity from $O(n^2 d)$ to $O(n r d)$ where $r$ is the feature dimension.

\subsubsection{Linformer}

The Linformer \cite{wang2020linformer} achieves linear complexity by projecting key and value matrices to a lower-dimensional space before computing attention. Based on the empirical observation that the attention matrix is approximately low-rank, Linformer applies learned linear projections to compress keys and values from length $n$ to a fixed dimension $k$, reducing the attention matrix from $n \times n$ to $n \times k$. This achieves 1.5$\times$ faster inference and 1.7$\times$ larger batch sizes while maintaining comparable performance.

\subsubsection{Kernel-Based Linear Attention and the Synthesizer}

Katharopoulos et al.\ demonstrated that transformers can be reformulated as recurrent neural networks when attention is linearized through kernel feature maps, achieving $O(n)$ complexity and up to 4,000$\times$ speedup on very long sequences. This formulation reveals a deep connection between transformers and RNNs. The Synthesizer challenged the necessity of query-key dot-product interaction by learning synthetic attention weights without token-token interactions, achieving comparable performance to standard transformers while being 60\% faster, suggesting that the specific content-based attention mechanism is less important than general position-aware information mixing.

\subsection{Hardware-Aware Attention}

\subsubsection{FlashAttention}

FlashAttention \cite{dao2022flashattention} achieves exact standard attention computation with dramatically improved practical performance by restructuring the computation to minimize memory accesses to high-bandwidth memory. The key insight is that the bottleneck in attention computation on modern GPUs is memory IO rather than floating-point operations. FlashAttention uses a tiling strategy that decomposes the attention computation into blocks fitting in fast on-chip SRAM, computing partial softmax results combined without ever materializing the full attention matrix in global memory.

FlashAttention reduces memory usage from $O(n^2)$ to $O(n)$ while computing exact attention, achieving 2--4$\times$ wall-clock speedup. FlashAttention-2 \cite{dao2023flashattention2} further improved performance through better parallelism, achieving near-optimal GPU utilization. FlashAttention has become the de facto standard attention implementation in production systems, demonstrating that hardware-aware algorithm design can provide improvements comparable to algorithmic approximations while maintaining exact computation.

\subsection{State-Space Models and Attention Alternatives}

\subsubsection{Structured State Spaces and Mamba}

The Structured State Space (S4) model \cite{gu2022efficiently} adapts continuous-time state-space models from control theory to sequence modeling, achieving $O(n \log n)$ complexity through convolution in the frequency domain. S4 parameterizes a linear dynamical system with specific structural constraints (the HiPPO initialization) that enable the model to capture long-range dependencies spanning thousands of positions. State-space models can be computed either as convolutions (efficient for training) or as recurrences (efficient for inference), providing optimal performance across computational settings.

Mamba \cite{gu2023mamba} extends the state-space framework with selective state spaces, where the state-space parameters are functions of the input, enabling content-dependent filtering. Mamba achieves linear-time sequence modeling with performance matching or exceeding transformers of equivalent size, while providing 5$\times$ faster inference throughput on long sequences. The selective state-space mechanism can be viewed as a content-dependent gating operation on a recurrent state, combining the efficiency of recurrent models with the content-sensitivity of attention.

\subsubsection{Hybrid Architectures}

Jamba \cite{lieber2024jamba} interleaves Mamba layers with transformer attention layers, using Mamba for efficient long-range modeling and attention for precise content-based retrieval. The hybrid architecture achieves performance superior to either pure approach and offers 3$\times$ throughput improvement over pure transformers on long sequences. RWKV \cite{peng2023rwkv} reinvents RNN-style computation using a linear attention variant with time-dependent decay, supporting both parallel training and efficient recurrent inference. RetNet \cite{sun2023retentive} proposes a retention mechanism supporting parallel, recurrent, and hybrid computation modes. These hybrid approaches represent a convergence toward architectures that dynamically trade off expressiveness and efficiency.

\subsection{Block-Sparse and Mixture-of-Attention Methods}

Beyond the foundational sparse patterns established by Longformer and BigBird, recent work has explored more flexible and adaptive sparse attention structures that move from fixed patterns to learned or dynamic sparsity.

\subsubsection{Block-Sparse Attention}

Block-sparse attention patterns divide the attention matrix into fixed-size blocks (typically 64--256 tokens), applying sparsity at the block level rather than the token level. This provides two advantages: better memory locality (blocks are contiguous in memory, enabling efficient GPU memory access) and more efficient hardware utilization (block operations map naturally to GPU tensor cores designed for matrix multiplication on fixed-size blocks). Block-sparse patterns have been implemented in libraries like Triton and integrated into FlashAttention, enabling practical deployment of sparse attention at production scale.

\subsubsection{Mixture of Block Attention}

The Mixture of Block Attention (MoBA) approach applies Top-K block selection, where each query block attends to only the most relevant key-value blocks as determined by a lightweight scoring mechanism. This approach bridges sparse attention and content-based retrieval: rather than using fixed patterns (local windows, global tokens), MoBA dynamically selects which parts of the context are most relevant for each query position. The scoring mechanism operates at the block level (comparing block-level summary representations), making it much cheaper than full token-level attention scoring while still capturing content-dependent relevance.

\subsubsection{Long-Context Generalization with Sparse Attention}

Recent work combines sparse attention mechanisms with length extrapolation techniques for improved long-context generalization. Hybrid approaches combining approximation with scaling consistently outperform single-method approaches, suggesting that the best long-context systems will use multiple efficiency techniques in concert. The Sparse Frontier meta-analysis provides a comprehensive analysis of sparse attention trade-offs, examining when sparsity helps and when full attention is necessary, and offering a decision framework for choosing sparse versus dense attention based on task type, sequence length, and hardware constraints. The key insight is that the optimal attention pattern is task-dependent: retrieval tasks benefit from global attention to key positions, reasoning tasks benefit from chains of local attention, and summarization tasks benefit from hierarchical attention that combines local processing with global aggregation.

\subsubsection{Efficient Attention Surveys and Systematization}

A comprehensive survey of efficient attention mechanisms for large language models organizes the landscape by mechanism type: sparse attention (fixed patterns, learned patterns, dynamic patterns), linear attention (kernel approximations, low-rank projections), hardware-optimized attention (tiling, memory-aware scheduling), and hybrid approaches (combining attention with alternative sequence mixing). The survey identifies a convergence toward adaptive methods that select attention patterns based on input content and task requirements, rather than applying fixed efficiency strategies. This adaptivity represents the attention-level analogue of the adaptive memory management advocated throughout this survey.

\subsection{Comparative Analysis and Design Principles}

Four broad conclusions emerge from the landscape of efficient attention architectures, each informing the design of context management systems.

First, hardware-aware exact attention (FlashAttention) has proven more practically impactful than algorithmic approximations for moderate context lengths (up to approximately 32K tokens), because it provides substantial speedup without quality degradation. The lesson is that implementation-level optimization---restructuring computation to match hardware characteristics---can provide improvements on par with algorithmic innovations while maintaining exact computation.

Second, for very long contexts (100K+ tokens), sparse or hybrid approaches become necessary as even FlashAttention's $O(n^2)$ time complexity becomes prohibitive. The quadratic wall is real: at 1 million tokens, even FlashAttention requires approximately $10^{12}$ operations per attention layer, necessitating either sparse patterns that reduce the number of computed attention scores or alternative architectures that avoid full pairwise attention entirely.

Third, the trend toward hybrid architectures combining attention with state-space or recurrent components suggests convergence on designs that use attention selectively for content-based retrieval while using more efficient mechanisms for general sequence processing. This design principle---using the most expensive mechanism (full attention) only where it is most needed (precise content-based retrieval)---recurs in the hierarchical memory architectures discussed in Section~\ref{sec:hierarchical}.

Fourth, linear attention and state-space models demonstrate that the specific content-based attention mechanism may be less important than general position-aware information mixing for many tasks. The Synthesizer's finding that learned synthetic attention weights achieve comparable performance to standard attention, while being 60\% faster, challenges the assumption that query-key dot-product interaction is essential, and suggests that future architectures may further reduce the role of full attention in favor of cheaper but sufficient alternatives.


% ============================================================
% SECTION 4: KV CACHE OPTIMIZATION
% ============================================================
\section{Key-Value Cache Optimization}
\label{sec:kvcache}

During autoregressive generation, large language models produce tokens one at a time, with each new token's computation depending on the key-value representations of all previous tokens. The key-value (KV) cache stores these representations to avoid redundant computation, but its memory requirements grow linearly with both sequence length and model size. For a model with $L$ layers, $h$ attention heads, and head dimension $d_h$, the KV cache for a sequence of length $n$ requires $2 \times L \times h \times d_h \times n$ elements, consuming tens of gigabytes for long sequences on large models. A typical 8-billion-parameter model may require loading 280 gigabytes of KV cache for long sequences---30$\times$ larger than the model parameters themselves---with inference efficiency dropping from approximately 70\% during training to roughly 10\% during generation due to memory-bandwidth limitations.

The KV cache optimization problem is fundamentally a memory management problem, and solutions draw on concepts from computer science (caching, paging, quantization) and information theory (rate-distortion trade-offs, importance sampling). This section organizes optimization strategies into four categories: eviction-based methods that reduce the number of cached tokens, quantization-based methods that reduce the precision of cached values, architectural innovations that reduce cache requirements through model design, and dynamic memory management systems that improve cache utilization through better resource allocation.

\subsection{Eviction-Based Cache Compression}

Eviction-based methods reduce cache size by selectively removing key-value pairs deemed least important for future generation, retaining only the most informative subset. The central challenge is determining which tokens are ``important''---a determination that ideally depends on future queries that are not yet known at eviction time. Different methods approximate importance using various signals, including cumulative attention scores, locality-sensitive hashing, and per-layer attention pattern analysis.

\subsubsection{Heavy-Hitter Oracle}

The Heavy-Hitter Oracle (H$_2$O) \cite{zhang2024h2o} identifies heavy-hitter tokens receiving consistently high attention scores and retains their KV pairs while evicting low-attention tokens. H$_2$O operates under the empirically validated assumption that a small subset of tokens dominates the attention distribution and remains relatively stable across generation steps. By maintaining heavy-hitter tokens alongside a sliding window of recent tokens, H$_2$O achieves up to 5$\times$ KV cache compression with minimal quality degradation. The method has become a widely-used baseline for KV cache compression research.

\subsubsection{Scissorhands and Persistence of Importance}

Scissorhands \cite{liu2023scissorhands} builds on the Persistence of Importance Hypothesis: tokens important at one generation step tend to remain important at subsequent steps. The method periodically scores all cached tokens by cumulative attention weight, evicts those below a threshold, and retains the rest. The scoring uses a variant of the Least Frequently Used (LFU) eviction policy from classical cache theory, combined with recency weighting to balance frequency and temporal proximity. This hypothesis holds across diverse models and tasks, achieving significant compression ratios with modest quality impact. However, the reliance on post-attention scoring means that attention must be computed before eviction decisions can be made, limiting the efficiency gains achievable through eviction alone.

\subsubsection{SnapKV and PyramidKV}

SnapKV \cite{li2024snapkv} observes that attention patterns during prompt processing reveal which tokens will be important during generation, enabling aggressive pre-generation compression. By observing attention patterns in the last few tokens of the prompt (the ``observation window''), SnapKV selects the tokens that will be most relevant for subsequent generation steps. This insight enables one-shot compression at the end of prompt processing rather than incremental eviction during generation. PyramidKV \cite{cai2024pyramidkv} extends this with a layer-aware approach based on the observation that different layers exhibit different attention concentration patterns (a pyramidal information funnel): early layers distribute attention broadly (requiring larger caches), while later layers concentrate attention on fewer tokens (requiring smaller caches). PyramidKV allocates cache budgets in a pyramidal pattern, with decreasing budgets for deeper layers, achieving better quality-efficiency trade-offs than uniform allocation.

\subsubsection{Adaptive KV Cache Compression}

FastGen \cite{ge2024model} uses the model's own attention patterns to determine which tokens to discard, profiling head-specific behaviors during a prefill phase and assigning each head an optimal eviction strategy from a portfolio of options: sliding window (for recency-focused heads), special token retention (for heads exhibiting attention sink behavior), heavy-hitter selection (for content-focused heads), and sparse top-k selection (for highly concentrated heads). This per-head customization achieves 40--50\% cache reduction with less than 1\% quality degradation and represents a principled approach where the model's internal behavior guides its own optimization.

\subsubsection{HashEvict: Pre-Attention Eviction}

HashEvict introduces locality-sensitive hashing (LSH) for pre-attention KV cache eviction, addressing a fundamental limitation of attention-based methods: they require computing attention scores before making eviction decisions, which partially defeats the purpose of reducing computation. HashEvict instead computes LSH hashes of query and key vectors and uses Hamming distances between hashes to approximate similarity, evicting tokens whose keys are dissimilar to the current query without computing full attention. This achieves 30--70\% KV cache compression on Llama~3 with 1.5--2$\times$ speedup over H$_2$O and Scissorhands during the prefill phase, with minimal quality degradation. The pre-attention approach is particularly valuable for latency-critical applications where reducing computation per step matters more than absolute compression ratio.

\subsubsection{Q-HITTER: Combined Eviction and Quantization}

Q-HITTER identifies a critical interaction between eviction and quantization: combining H$_2$O's token selection with low-bit quantization degrades performance because quantization distorts the attention patterns used for selection, causing the wrong tokens to be retained. Q-HITTER proposes a dual-metric scoring system combining attention importance with quantization friendliness---tokens with distributions amenable to low-bit representation score higher. This combined approach achieves up to 20$\times$ memory reduction with throughput improvements of 33$\times$ over Hugging Face Accelerate and 7$\times$ over DeepSpeed, demonstrating that jointly optimizing eviction and quantization outperforms applying either technique independently.

\subsection{Token Merging and Reversible Compression}

Beyond eviction (which permanently discards tokens), merging approaches combine multiple tokens into compressed representations, and reversible methods enable reconstruction of compressed tokens when needed.

KVReviver uses a sketch algorithm for KV cache compression that allows reconstruction of compressed tokens on demand, trading memory for computation by storing compressed sketches and reconstructing full representations when needed. RocketKV employs a two-stage approach combining permanent eviction (for clearly unimportant tokens) with dynamic per-batch selection (optimizing which remaining tokens to use for each generation step). MiniCache \cite{jia2024minicache} merges KV cache across depth (layers), consolidating states every $N$ layers based on the observation that adjacent layers often produce highly correlated representations. This depth-wise merging reduces cache size significantly while preserving the information content distributed across layers.

\subsection{Quantization-Based Cache Compression}

Quantization reduces the precision of cached key-value representations, reducing memory per cached token rather than the number of cached tokens. This approach is orthogonal to eviction and can be combined with it for multiplicative compression gains.

KIVI \cite{liu2024kivi} demonstrates that KV cache values can be quantized to 2-bit precision with minimal quality impact using asymmetric per-channel key quantization and per-token value quantization, achieving approximately 8$\times$ memory reduction. The key insight is that keys and values have different quantization characteristics: keys exhibit per-channel outliers (requiring per-channel scaling), while values exhibit per-token outliers (requiring per-token scaling). KVQuant \cite{hooper2024kvquant} targets 10-million-token contexts through aggressive quantization with residual outlier handling, achieving near-lossless quality at 2--4 bit precision across models up to 70B parameters. Custom CUDA kernels achieve 1.7$\times$ speedups for quantized key and value matrix-vector multiplications at 4-bit precision on LLaMA-7B. GEAR \cite{kang2024gear} combines quantization with low-rank approximation of quantization residuals, enabling more aggressive base quantization without quality loss by capturing quantization error in a separate low-rank representation.

\subsection{Architectural Cache Optimization}

Grouped Query Attention \cite{ainslie2023gqa} reduces cache size at the architectural level by sharing KV heads across groups of query heads, achieving 4$\times$ or greater reduction. For Llama~3 8B, GQA reduces KV cache from 280 gigabytes to 69 gigabytes---a 4$\times$ reduction that fundamentally changes deployment economics. Cross-Layer Attention \cite{brandon2024reducing} observes that KV representations in adjacent layers are often highly correlated (correlation coefficients exceeding 0.9 for adjacent layers) and proposes sharing caches across layers, achieving further 2$\times$ reduction. The combination of GQA and cross-layer sharing can reduce cache memory by 8$\times$ or more without changing the attention algorithm, representing a pure architectural optimization.

\subsection{Dynamic Memory Management}

\subsubsection{PagedAttention and vLLM}

PagedAttention \cite{kwon2023efficient} applies virtual memory concepts from operating systems to KV cache management. Traditional KV cache implementations pre-allocate contiguous memory for the maximum possible sequence length, leading to massive internal fragmentation---often 60--80\% of allocated memory is wasted. PagedAttention divides each request's KV cache into fixed-size blocks (analogous to OS memory pages) allocated on demand from a shared pool, with a block table mapping logical cache positions to physical memory locations. This achieves near-zero memory waste, enables efficient memory sharing between requests with common prefixes (analogous to copy-on-write in OS memory management), and supports dynamic batching that maximizes GPU utilization. The vLLM serving system built on PagedAttention has become the standard framework for LLM inference in production, achieving 2--4$\times$ higher throughput than HuggingFace Transformers and 2.2$\times$ higher than the Text Generation Inference server.

\subsubsection{vAttention}

vAttention, published at ASPLOS~2025, proposes leveraging operating system demand paging instead of application-level paging for KV cache management. While PagedAttention implements its own virtual memory system at the application level (requiring custom attention kernels), vAttention allocates KV cache in contiguous virtual memory and relies on the OS kernel to manage physical page allocation through demand paging. Standard attention kernels work without modification because the OS transparently handles physical memory through page faults. This achieves 1.97$\times$ faster token generation than vLLM and 3.92$\times$ faster prompt processing than FlashAttention variants, with dramatically simplified implementation. The success of vAttention demonstrates that well-designed OS abstractions can outperform application-level reimplementations of the same concepts.

\subsubsection{Entropy-Guided Budget Allocation}

Entropy-guided KV cache budget allocation uses Shannon entropy of each layer's attention distribution to dynamically allocate cache budgets. The entropy $H = -\sum_i p_i \log p_i$ of the attention distribution measures how dispersed versus concentrated attention is: high-entropy layers distribute attention broadly across many tokens and require larger caches to retain the diverse set of attended tokens, while low-entropy layers concentrate attention on few tokens and function with smaller caches. This principled allocation based on information-theoretic criteria achieves memory reductions and 46.6\% decoding time reduction with preserved quality.

\subsubsection{Prompt Caching}

Prompt caching addresses redundant computation when many requests share common prompt components---a pattern ubiquitous in production deployments where system prompts, few-shot examples, or document contexts are reused across requests. By pre-computing and caching KV states for reusable prompt segments, subsequent requests compute attention only for novel portions. This achieves 8--60$\times$ improvement in time-to-first-token latency and has been adopted by major API providers including OpenAI, Anthropic, and Google. The effectiveness of prompt caching depends on the prefix-sharing structure of incoming requests, with the highest gains in applications (such as customer service bots or code assistants) where a substantial system prompt is shared across all interactions.

\subsection{Eviction Strategy Selection}

The diversity of KV cache optimization methods raises the practical question of which strategy to select for a given deployment scenario. For streaming applications requiring constant memory, the combination of attention sinks with a sliding window (StreamingLLM) provides simplicity and reliability. For batch inference workloads, H$_2$O or Scissorhands offers a good quality-efficiency trade-off. For long contexts where quality is paramount, Q-HITTER's combined eviction and quantization provides the best overall efficiency. For latency-critical applications, HashEvict's pre-attention computation avoids the overhead of full attention scoring. In practice, the optimal strategy depends on the specific hardware, batch size, sequence length distribution, and quality requirements of the deployment, and many production systems implement adaptive selection among multiple strategies based on runtime conditions.


% ============================================================
% SECTION 5: CONTEXT COMPRESSION
% ============================================================
\section{Context Compression and Information Distillation}
\label{sec:compression}

Context compression reduces the number of tokens that need to be processed by distilling essential information from a long context into a shorter representation. Unlike the KV cache optimization methods of Section~\ref{sec:kvcache}, which operate on intermediate representations within the model, context compression methods operate on the input itself, producing a shorter prompt that preserves the essential information content of the original. This section surveys compression methods ranging from learned token compression to prompt distillation frameworks, token pruning approaches, and the information-theoretic foundations that establish fundamental limits on compression achievability.

\subsection{Learned Token Compression}

Learned compression methods train models to map long token sequences into shorter representations that can be processed by a standard language model, achieving compression through learned encodings rather than heuristic selection.

\subsubsection{Gist Tokens}

Mu et al.\ \cite{mu2024learning} introduced gisting, which compresses input prompts into gist tokens by introducing a representational bottleneck through attention mask modification. The implementation requires only new vocabulary token embeddings and a modified attention mask that forces the model to route all information through the gist token positions. The model learns to condense information into dedicated gist token activations in a single forward pass, generating a context-specific ``tuning prefix'' that captures the essential content of the compressed prompt. However, performance degrades rapidly beyond contexts of a few hundred tokens, as the fixed number of gist tokens becomes insufficient to represent the information content of longer inputs. This limitation was characterized in subsequent work demonstrating that gisting does not effectively scale even to contexts of hundreds of tokens, making it unsuitable for the long-context scenarios that motivate modern context management research.

\subsubsection{AutoCompressors}

AutoCompressors \cite{chevalier2023adapting} extend gisting to longer contexts by training language models to recursively compress context into compact summary vectors. The model processes input in segments, compressing each segment into a fixed number of summary tokens that are prepended to the next segment's input. This recursive structure enables processing of arbitrarily long documents within a bounded memory footprint: a document of $n$ tokens is processed in $n/s$ segments of size $s$, with each segment producing $k$ summary tokens, maintaining a constant memory overhead of $k$ tokens regardless of document length. AutoCompressors achieve significant perplexity improvements over truncation baselines, demonstrating that recursive compression preserves substantially more information than simply discarding distant context. The recursive architecture directly parallels the recursive summarization approaches used in multi-turn dialogue systems discussed in Section~\ref{sec:multiturn}.

\subsubsection{In-Context Autoencoder}

The In-Context Autoencoder (ICAE) \cite{ge2024icae} trains an autoencoder within the context window, with an encoder compressing long context into memory tokens and a decoder reconstructing relevant information during generation. Unlike gisting, which uses a single forward pass, ICAE explicitly trains for reconstruction quality through an autoencoder objective that ensures the compressed representation retains retrievable information. ICAE achieves compression ratios of 4$\times$ or higher while maintaining competitive generation quality, suggesting that long contexts contain substantial exploitable redundancy. The autoencoder formulation provides a principled training objective---minimize reconstruction error---that connects context compression to the broader literature on learned data compression.

\subsection{Prompt Compression Frameworks}

Prompt compression frameworks operate at inference time, reducing the length of input prompts through selective token removal or importance-weighted filtering. These methods do not require model modification and can be applied to any language model as a preprocessing step.

\subsubsection{LLMLingua}

LLMLingua \cite{jiang2023llmlingua} uses a small language model to estimate token-level information content and selectively removes low-information tokens. The compression procedure operates in three stages. First, a budget controller determines the overall compression ratio based on the length and content type of the input. Second, a demonstration-level compressor selects the most informative few-shot examples when demonstrations are present. Third, a token-level iterative compressor uses perplexity scores from the small model to identify and remove predictable tokens, retaining tokens with high perplexity (high information content) while discarding those the model could predict from surrounding context. LLMLingua achieves up to 20$\times$ compression with minimal performance degradation on reasoning and question-answering tasks, demonstrating that the majority of tokens in typical prompts are informationally redundant from the model's perspective.

\subsubsection{LongLLMLingua}

LongLLMLingua \cite{jiang2024longllmlingua} extends prompt compression to long-context scenarios with three key innovations. First, question-aware compression conditions token importance on the downstream query, ensuring that tokens relevant to the specific question are retained even if they would be considered low-information in isolation. Second, document reordering places the most relevant retrieved documents at the beginning and end of the compressed prompt, explicitly mitigating the lost-in-the-middle effect described in Section~\ref{sec:attention}. Third, dynamic compression ratios apply different levels of compression to different parts of the context based on their estimated relevance. LongLLMLingua achieves simultaneous improvements in both efficiency and quality, as compression removes noisy context that would otherwise degrade performance---demonstrating that for long contexts, less can genuinely be more.

\subsubsection{Selective Context}

Selective Context \cite{li2023compressing} evaluates the self-information of each lexical unit under a causal language model and removes units with low self-information, preserving informative surprise content while discarding predictable filler. The self-information $I(x) = -\log p(x | \text{context})$ provides a principled, information-theoretic criterion for token importance: tokens that the model can predict well carry little information and can be safely removed. The approach is computationally lightweight and applicable as a preprocessing step to any language model without modification.

\subsection{Token Pruning and Dynamic Computation}

Token pruning approaches reduce computation during inference by selectively skipping tokens or attention heads, achieving efficiency without permanently removing tokens from the input.

\subsubsection{LazyLLM: Dynamic Token Pruning}

LazyLLM, developed by Apple Machine Learning Research, enables language models to dynamically select different subsets of tokens from context in different generation steps. Rather than computing attention over all tokens for every generated token, LazyLLM defers computation of less important tokens to later steps when they may become relevant. Token importance is measured using attention scores from prior transformer layers: if a token receives low attention in early layers, its KV computation is deferred in subsequent layers. This lazy evaluation strategy reduces computation for long sequences while maintaining token availability---deferred tokens are not permanently evicted but remain accessible when needed. The key distinction from eviction-based methods is reversibility: LazyLLM can always recompute deferred tokens if they become relevant, avoiding the information loss inherent in permanent eviction.

\subsubsection{Attention Head Pruning}

Complementing token-level pruning, attention head pruning removes unnecessary heads from the multi-head attention mechanism. Research has demonstrated that only a small subset of attention heads is important for downstream tasks, with important heads exhibiting interpretable functions such as attending to adjacent words or tracking specific syntactic relations. Layer-wise pruning analyzes head importance per layer and applies layer-specific pruning masks. Single-shot meta-pruning compresses pre-trained transformers before fine-tuning, adaptively pruning heads for different downstream tasks. A* search-based pruning provides strict accuracy guarantees while eliminating redundant heads through optimal search. These structural pruning methods can reduce parameters by 30--60\% while maintaining or even improving task performance, with the additional benefit of reduced KV cache requirements for pruned heads.

\subsubsection{Knowledge Distillation for Context Efficiency}

Knowledge distillation transfers the capabilities of large ``teacher'' models to smaller ``student'' models, enabling more efficient context processing. The Minitron approach combines weight pruning with knowledge distillation to reduce training costs of model families, producing smaller and faster models that maintain high performance. The distillation process involves output distillation (matching final predictions), feature distillation (matching intermediate representations), and relation distillation (matching relationships between data samples). For context management, distillation is particularly valuable because smaller student models have proportionally smaller KV caches, enabling longer effective context windows on the same hardware.

\subsection{Compressive Transformers}

The Compressive Transformer \cite{rae2020compressive} integrates compression into the model architecture, maintaining a two-level memory hierarchy: fine-grained memory of recent tokens and coarser compressed memory of older tokens. As tokens age out of fine-grained memory, they are compressed through learned compression functions---which can include mean pooling, 1D convolution, dilated convolution, or attention-based compression---and stored in compressed memory. This architectural approach mirrors the human memory system's transition from detailed short-term representations to schematic long-term representations, demonstrating that graceful information degradation is more effective than hard truncation. The Compressive Transformer anticipates the hierarchical memory architectures discussed in Section~\ref{sec:hierarchical} and provides early evidence that multi-resolution memory representations improve long-context performance.

\subsection{Information-Theoretic Perspectives on Compression}

The effectiveness of context compression can be understood through information theory, which provides both motivation for compression and fundamental limits on what compression can achieve.

\subsubsection{Entropy and Redundancy in Natural Language}

Natural language exhibits substantial redundancy: the entropy rate of English text is estimated at 1.0--1.5 bits per character, far below the theoretical maximum of approximately 4.7 bits per character for a 26-letter alphabet plus space. This redundancy implies that a significant fraction of tokens can be removed without meaningful information loss. Shannon's source coding theorem establishes that no lossless compression scheme can achieve a compression ratio below the entropy rate, providing a hard floor on achievable compression. For typical English text, this theoretical limit suggests compression ratios of approximately 3$\times$ are achievable without information loss---but practical methods like LLMLingua achieve 20$\times$ compression by accepting lossy compression of task-irrelevant content.

\subsubsection{Rate-Distortion Framework}

The rate-distortion framework from information theory formalizes the trade-off between compression ratio and reconstruction quality. For context compression in LLMs, the ``distortion'' is the degradation in downstream task performance caused by compression, and the ``rate'' is the number of tokens retained after compression. The rate-distortion function $R(D)$ defines the minimum rate (compressed length) needed to achieve distortion at most $D$ (acceptable quality loss), establishing hard limits on useful compression ratios for a given quality target. Research on information-theoretic limits has shown that with a frozen LLM, there exists a hard limit on compression achievable through token dropping or reordering, determined by the information content of the context relative to the downstream task.

\subsubsection{LLMs as Compression Engines}

A deep connection exists between language modeling and data compression. LLMs output next-token probability distributions that can be used for near-optimal compression via arithmetic coding: tokens that the model predicts with high probability require few bits to encode, while surprising tokens require more bits. This dual nature---LLMs as both language models and compression engines---means that compression quality is directly related to language modeling quality. Better language models achieve better compression, and the compression ratio achieved by a language model on a given text provides a measure of how well the model understands that text. This perspective connects context compression to the broader question of what LLMs ``know'' about their inputs and establishes that compression-based methods are not merely engineering heuristics but are grounded in fundamental information-theoretic principles.

\subsubsection{Compression-Performance Trade-offs}

The relationship between compression ratio and downstream performance is fundamentally non-linear. At low compression ratios (retaining 80--90\% of tokens), performance is nearly identical to the uncompressed baseline for most tasks. At moderate ratios (retaining 30--50\%), performance degrades gracefully, with task-aware compression methods (LongLLMLingua) significantly outperforming task-agnostic methods (random dropping). At extreme ratios (retaining less than 10\%), performance degrades sharply for information-dense tasks (question answering, fact verification) but may remain acceptable for tasks requiring only gist understanding (summarization, sentiment analysis). This task-dependent degradation pattern echoes selective encoding in human memory, where goal-relevant information receives preferential encoding while task-irrelevant details are forgotten without consequence.


% ============================================================
% SECTION 6: MEMORY-AUGMENTED ARCHITECTURES
% ============================================================
\section{Memory-Augmented Architectures}
\label{sec:memory}

Memory-augmented architectures supplement the model's internal context with external memory systems that can store, retrieve, and manage information independently of the fixed context window. This paradigm shift---from extending the model to augmenting it---has produced some of the most practically impactful advances in context management. Rather than making the model's context window larger, memory-augmented approaches give the model access to effectively unlimited external information through retrieval mechanisms, creating a hybrid system where the context window serves as working memory while external storage serves as long-term memory. This section traces the evolution from classical neural memory through modern retrieval-augmented generation to operating-system-inspired memory management.

\subsection{Classical Neural Memory}

The concept of augmenting neural networks with external memory predates the transformer era and provides important conceptual foundations for modern systems. Graves et al.\ \cite{graves2014neural} introduced Neural Turing Machines (NTMs), coupling a neural network controller with an external memory matrix accessed through differentiable read and write operations. The NTM controller learns to use memory through content-based addressing (retrieving by similarity to a query) and location-based addressing (sequential or relative position access), establishing a vocabulary of memory operations that recurs in modern systems. The Differentiable Neural Computer (DNC) extended NTMs with dynamic memory allocation and temporal link tracking, enabling more sophisticated memory management.

Memory Networks \cite{weston2015memory} simplified external memory with embedding-based retrieval over a memory bank, framing memory access as an attention operation over stored memories. End-to-End Memory Networks \cite{sukhbaatar2015end} made the entire pipeline differentiable, enabling end-to-end training of reading and writing operations. These architectures established three foundational concepts that directly inform modern context management: external storage decoupled from the model's parameters, content-based retrieval using learned similarity functions, and iterative refinement through multiple memory access steps (``hops'').

The Memorizing Transformer \cite{wu2022memorizing} bridged classical neural memory and modern transformers by augmenting a standard transformer with a non-differentiable approximate nearest-neighbor memory of past key-value pairs. The model retrieves the most relevant past tokens using $k$-nearest-neighbor search (specifically FAISS indexing) and incorporates them into the attention computation through an additional cross-attention layer. This enables effective context to extend to hundreds of thousands of tokens without modifying the base model's architecture or training procedure, demonstrating that external memory can provide dramatic context extension with minimal architectural changes.

\subsection{Retrieval-Augmented Generation}

Retrieval-Augmented Generation (RAG) \cite{lewis2020retrieval} represents the most practically impactful memory augmentation paradigm, combining a parametric language model with a non-parametric retrieval system that fetches relevant documents from an external corpus at inference time. The foundational RAG paper by Lewis et al.\ demonstrated two formulations: RAG-Sequence (using the same retrieved passages for the entire generated sequence) and RAG-Token (potentially using different passages for each generated token). RAG achieved state-of-the-art results on Natural Questions, TriviaQA, and WebQuestions, establishing retrieval augmentation as a standard approach for knowledge-intensive tasks.

\subsubsection{Dense Retrieval Foundations}

Dense Passage Retrieval (DPR) \cite{karpukhin2020dense} replaced traditional sparse retrieval (BM25) with learned dense vector representations, using a dual-encoder architecture with separate BERT-based encoders for queries and passages. DPR achieves 9--19\% higher top-20 retrieval accuracy than BM25 through learned semantic matching that overcomes the lexical mismatch problem inherent in sparse methods. FAISS indexing enables efficient similarity search over millions of passages. Sentence-BERT \cite{reimers2019sentencebert} further improved dense retrieval with efficient sentence-level embeddings optimized through contrastive learning. ColBERT \cite{khattab2020colbert} introduces late interaction retrieval with per-token embeddings for queries and documents, preserving fine-grained matching information while remaining efficient through MaxSim operations over precomputed token embeddings. ColBERTv2 \cite{santhanam2022colbertv2} improved efficiency through residual compression of document embeddings, reducing storage requirements by 6--10$\times$.

\subsubsection{RETRO}

RETRO \cite{borgeaud2022improving} integrates retrieval into pre-training itself, retrieving relevant passages from a 2-trillion-token database during both training and inference. Unlike standard RAG where retrieval is applied only at inference, RETRO trains the model to effectively condition on retrieved context through a chunked cross-attention mechanism where the input is divided into chunks and each chunk attends to its nearest-neighbor retrieved passages. RETRO achieves performance comparable to models 25$\times$ larger, suggesting that much parametric knowledge could be more efficiently stored in external retrieval corpora rather than encoded in model parameters. This finding has profound implications for the scaling debate: rather than training ever-larger models, it may be more efficient to combine moderate-sized models with extensive retrieval corpora.

\subsubsection{Self-RAG}

Self-RAG \cite{asai2024selfrag} trains the language model to generate special reflection tokens that control retrieval and generation behavior. A Retrieval token decides whether external retrieval is needed for the current generation step (avoiding unnecessary retrieval when the model has sufficient parametric knowledge). A Relevance token evaluates whether each retrieved passage is relevant to the query. A Support token assesses whether the generation is supported by the retrieved evidence. A Critique token provides overall quality assessment. This self-reflective approach reduces hallucination by avoiding unsupported claims, improves efficiency by skipping unnecessary retrieval, and provides interpretable decision traces. Self-RAG models at 7B and 13B parameters outperform ChatGPT and Llama~2-Chat on open-domain QA, fact verification, and reasoning tasks while using retrieval selectively rather than for every query.

\subsubsection{Corrective RAG}

CRAG \cite{yan2024corrective} introduces a lightweight retrieval evaluator that computes confidence scores for retrieved documents, addressing the fundamental weakness of standard RAG: retrieved documents may be irrelevant, outdated, or misleading. When the retrieval evaluator's confidence falls below a threshold, corrective actions are triggered. These include web search for supplementary information, query expansion to reformulate the retrieval query, and query decomposition to break complex queries into sub-queries that can be retrieved independently. CRAG significantly reduces hallucinations caused by poor retrieval and is designed as a plug-and-play module compatible with any RAG pipeline, making it straightforward to retrofit onto existing systems.

\subsubsection{GraphRAG}

GraphRAG constructs knowledge graphs from source documents through entity and relationship extraction, organizes them into hierarchical communities using graph clustering algorithms, and generates summaries at each community level. This structured representation enables two complementary search modes: local search retrieves information about specific entities by traversing graph neighborhoods, combining entity descriptions, relationship context, and community summaries; global search aggregates community summaries for broad questions requiring synthesis across the entire corpus. GraphRAG achieves up to 35\% improvement in precision over vector-only retrieval on complex multi-entity queries and enables multi-hop reasoning through explicit relationship traversal---a capability that flat vector-based RAG systems struggle to achieve because they lack structural information about how entities relate.

\subsubsection{Advanced RAG Variants}

The RAG paradigm has proliferated into numerous specialized variants addressing specific limitations. Naive RAG (retrieve-then-generate) represents the simplest pipeline. Advanced RAG incorporates iterative retrieval with query refinement, where initial retrieval informs query reformulation for improved subsequent retrieval. Modular RAG separates retrieval, reranking, filtering, and generation into composable modules that can be independently optimized. Agent-Based RAG gives the language model agency to decide when, what, and how to retrieve through tool-use interfaces, representing the convergence of RAG with agentic architectures. Hybrid retrieval combines dense and sparse methods, using BM25 for lexical matching and dense retrieval for semantic matching, with learned fusion weights that adapt to query characteristics.

\subsection{Operating-System-Inspired Memory Management}

\subsubsection{MemGPT}

MemGPT \cite{packer2023memgpt} draws an explicit analogy between OS virtual memory management and LLM context management, treating the LLM's context window as ``main memory'' (fast but limited) and external databases as ``disk storage'' (slow but unlimited). The architecture implements a multi-tier memory system. Tier~1 is the main context---the LLM's actual context window---containing core memory blocks for persona (agent identity and goals), human (user characteristics and relationship), and scratch (temporary working memory for the current task), typically consuming 1--2K tokens. Tier~2 is external storage containing the full document corpus, conversation history, and metadata, accessible through semantic search. Tier~3 consists of summaries and indexes that guide retrieval from Tier~2.

The critical innovation is that the LLM itself manages memory through explicit function calls---reading from external storage, writing to memory, consolidating information, and evicting less important content. Memory management is driven by an interrupt model analogous to OS context switches: when the context window approaches capacity, the system pauses generation, performs memory management operations (eviction, summarization, retrieval), and resumes generation. This self-directed memory management enables coherent conversations spanning thousands of turns, preference tracking across sessions, and adaptive information retrieval. MemGPT can analyze documents 10--20$\times$ larger than the context window with quality comparable to the base model on shorter documents. The system has been commercialized as the Letta platform, demonstrating practical viability of the OS-inspired memory paradigm.

\subsubsection{Infini-Attention}

Infini-Attention \cite{munkhdalai2024leave} integrates compressive memory directly into the attention mechanism, maintaining accumulated information from all past tokens using a linear attention-like update rule. The compressive memory is updated through an associative memory operation that maps keys to values, accumulating information without storing individual key-value pairs. A gating mechanism controls the balance between local attention (standard softmax attention over the current segment) and retrieved memory (information from the compressive store), enabling the model to adaptively choose between recent detailed information and compressed historical information. Infini-Attention enables infinite context within bounded memory, processing sequences of arbitrary length with fixed memory per layer.

\subsubsection{Episodic Memory Systems}

EM-LLM \cite{hu2025emllm} applies event segmentation theory from cognitive psychology to LLM memory management. The system segments incoming token streams into discrete episodes based on surprise signals---sudden changes in the model's internal representations that indicate event boundaries, analogous to human perception of event transitions. Episodes are stored in an external buffer with embeddings that enable similarity-based retrieval. When processing a new query, the system retrieves the most relevant episodes and incorporates them into the context, providing targeted historical information without the cost of maintaining the entire context. This episodic structure provides natural organization of information by event rather than by arbitrary token boundaries.

Think-in-Memory \cite{liu2023thinkinmemory} extends episodic memory with explicit thinking operations that mirror cognitive deliberation. The system retrieves relevant memories, reasons about how to integrate retrieved information with new input (the ``thinking'' step), generates a response, and stores updated understanding back to memory. This explicit reasoning step addresses a limitation of standard RAG: simply concatenating retrieved passages with the query does not guarantee effective integration. By adding an intermediate reasoning step, Think-in-Memory achieves better information synthesis, particularly for queries requiring multi-step reasoning across retrieved memories.


% ============================================================
% SECTION 7: HIERARCHICAL HUMAN-INSPIRED MEMORY
% ============================================================
\section{Hierarchical Human-Inspired Memory Architectures}
\label{sec:hierarchical}

The preceding sections surveyed technical approaches to context management primarily motivated by computational constraints. This section---the core of the present survey---argues that cognitive science provides essential design principles for context management, and that the most capable long-context architectures draw explicitly on theories of human memory organization. The central thesis is that human memory is not a single monolithic system but a collection of specialized, interacting subsystems organized hierarchically by temporal scale, abstraction level, and functional role. The most effective LLM memory architectures mirror this multi-system organization.

\subsection{Foundations in Cognitive Psychology}

\subsubsection{The Atkinson-Shiffrin Model}

The Atkinson-Shiffrin model \cite{atkinson1968human}, often called the modal model of memory, proposed three distinct memory stores through which information flows sequentially. The sensory register holds a brief, high-fidelity representation of incoming stimuli lasting milliseconds to seconds. Information receiving attention transfers to short-term memory (STM), a limited-capacity store holding approximately seven items (plus or minus two, as characterized by Miller \cite{miller1956magical}) for 15--30 seconds without rehearsal. Through rehearsal, information in STM may be encoded into long-term memory (LTM), a theoretically unlimited-capacity store.

The Atkinson-Shiffrin model introduced concepts that directly inform LLM memory design. Capacity-limited short-term storage maps to the transformer's finite context window. The distinction between maintenance rehearsal (sustaining items without promoting long-term encoding) and elaborative rehearsal (creating meaningful associations for long-term storage) parallels the difference between retaining tokens in the context window and actively compressing or summarizing context for persistent storage. The emphasis on control processes---voluntary strategies for encoding, retrieval, and rehearsal---anticipates self-directed memory management in systems like MemGPT. Neath and Surprenant \cite{neath2019fifty} reviewed fifty years of subsequent research, noting that while the strict three-store architecture has been refined, its core insight---multiple interacting systems with different characteristics---remains foundational.

\subsubsection{Baddeley's Working Memory Model}

Baddeley and Hitch \cite{baddeley1974working} replaced the unitary short-term store with a multi-component working memory system: a central executive functioning as an attentional control system, a phonological loop maintaining verbal information through rehearsal, a visuospatial sketchpad handling spatial information, and an episodic buffer \cite{baddeley2000episodic} integrating information from multiple sources into coherent episodes.

This architecture has direct implications for LLM design. The central executive parallels the language model itself in systems like MemGPT and CoALA, where the LLM decides what to attend to, retrieve, and store. Domain-specific slave systems suggest that different information types might benefit from specialized storage mechanisms. The episodic buffer's integration role is mirrored in RAG systems combining retrieved documents with parametric knowledge and conversational context. The working memory capacity limit of approximately four chunks is relevant: the effective information capacity of LLM context windows may be far smaller than nominal token counts, as the lost-in-the-middle effect demonstrates. This parallel suggests LLM architectures might benefit from explicit chunking and hierarchical organization rather than flat token sequences.

\subsubsection{Tulving's Memory Taxonomy}

Tulving \cite{tulving1972episodic, tulving1985memory} distinguished episodic memory (memory for personally experienced events anchored in spatial-temporal context) from semantic memory (memory for general knowledge abstracted from specific experiences). This distinction maps onto two fundamentally different types of information LLM systems must manage. Semantic memory corresponds to knowledge encoded in model parameters during pre-training: factual knowledge, linguistic patterns, and conceptual relationships accessed without reference to specific training examples. Episodic memory corresponds to information from specific interactions, conversations, or documents that must be stored and retrieved with contextual associations intact.

Current LLMs have sophisticated semantic memory but lack native episodic memory---they cannot intrinsically recall specific past interactions. This asymmetry has motivated explicit episodic memory systems: Generative Agents \cite{park2023generative} implement episodic memory through timestamped memory streams, MemoryBank \cite{zhong2024memorybank} implements decaying episodic memories, EM-LLM \cite{hu2025emllm} applies event segmentation theory, and Sun et al.\ \cite{sun2025episodic} explicitly argue that episodic memory is the missing piece for long-term agents. Ning et al.\ \cite{ning2025episodic} synthesized these parallels in Trends in Cognitive Sciences, arguing that successful architectures must implement temporal context, spatial context, self-referential encoding, and associative retrieval cues.

\subsubsection{Ebbinghaus and the Forgetting Curve}

Ebbinghaus \cite{ebbinghaus1885memory} established that memory retention decays exponentially: $R(t) = e^{-t/S}$, where $R$ is retention, $t$ is time, and $S$ is memory stability. Murre and Dros \cite{murre2015replication} replicated and confirmed these findings with modern methods. Two aspects are particularly relevant to LLM memory design. First, forgetting is adaptive: by allowing irrelevant information to decay, the system reduces interference and maintains retrieval efficiency---informing cache eviction and memory management strategies. Second, the spacing effect---distributed learning episodes produce stronger memories than concentrated ones---suggests that periodic retrieval and re-encoding of important memories can counteract decay, as implemented in MemoryBank's forgetting-curve-based management.

\subsubsection{McClelland's Complementary Learning Systems}

McClelland, McNaughton, and O'Reilly \cite{mcclelland1995complementary} proposed two complementary learning systems: the hippocampus, rapidly encoding specific experiences with minimal interference, and the neocortex, gradually extracting statistical regularities through slow, interleaved learning. Memory consolidation---gradual transfer from hippocampal to neocortical representations, often during sleep---bridges these systems.

This theory provides a powerful framework for LLM memory architecture. Model parameters, trained through gradient descent over billions of examples, function as the neocortical system: encoding general knowledge through slow, interleaved learning that extracts statistical regularities across the training corpus. The slow learning rate and massive dataset ensure that individual examples do not catastrophically interfere with previously learned knowledge, just as neocortical learning preserves existing knowledge through gradual integration. External memory systems (RAG databases, episodic buffers, conversation logs) function as the hippocampal system: rapidly storing specific experiences without parameter modification, through write operations that take effect immediately without requiring gradient updates. The hippocampal system's rapid encoding is essential because waiting for slow neocortical learning to incorporate new information would render the system unable to use recent experience.

Memory consolidation---the gradual transfer from hippocampal to neocortical representations---corresponds to fine-tuning, knowledge distillation, or summarization of episodic memories into semantic representations. In biological systems, consolidation occurs primarily during sleep through hippocampal replay coordinated with neocortical oscillations. The computational analogue involves periodic processing of accumulated episodic memories, potentially through parameter-efficient fine-tuning (LoRA adapters), summarization into persistent knowledge stores, or re-indexing of episodic memories for more efficient retrieval. The CLS framework explains why purely parametric and purely retrieval-based approaches are each insufficient: effective management requires both rapid hippocampal-like storage for immediate use and gradual neocortical-like integration for long-term knowledge accumulation.

The CLS framework also illuminates the phenomenon of catastrophic forgetting (discussed in Section~\ref{subsec:catastrophic}): when fine-tuning a model on new data, the rapid weight updates resemble hippocampal rather than neocortical learning, overwriting previously stored knowledge. This explains why continual learning requires special mechanisms (elastic weight consolidation, experience replay) that simulate the slow, interleaved learning characteristic of the neocortical system.

\subsubsection{Memory Consolidation and Sleep}

During slow-wave sleep, the hippocampus replays recently encoded experiences, coordinated with neocortical oscillations through spindle-ripple coupling \cite{diekelmann2010memory, carr2011hippocampal, foster2017replay}. Gonz{\'a}lez-Rueda et al.\ \cite{gonzalez2022model} modeled autonomous hippocampal-neocortical interactions during sleep, demonstrating that coordinated replay is essential for effective consolidation.

This has inspired computational analogues. Sleep Replay Consolidation (SRC) implements an offline phase where stored latent representations are replayed using Hebbian-like learning, strengthening important connections and reducing catastrophic forgetting. Experiments demonstrate that sleep consolidation preserves Task~1 knowledge at 86\% accuracy versus 18\% without sleep while learning Task~2. NeuroDream \cite{tutuncuoglu2024neurodream} extends this with an explicit dream phase using a learned dynamics model, generating synthetic experiences that consolidate learning without original training data. NeuroDream reduces forgetting from 40\% to 15\% at approximately 20\% increased training cost.

These approaches suggest LLM memory systems might benefit from periodic offline consolidation---scheduled intervals for reviewing, reorganizing, and consolidating stored memories through summarization, re-indexing, or adapter fine-tuning on accumulated experience.

\subsection{The CoALA Framework}

The Cognitive Architectures for Language Agents (CoALA) framework \cite{sumers2024cognitive} provides the most systematic taxonomy for LLM agent memory and action components, drawing on decades of cognitive architecture research including Soar \cite{laird2012soar}, ACT-R \cite{anderson2004integrated}, and the broader production system tradition. CoALA decomposes agent cognition into three interacting systems: memory (storage, organization, and retrieval of information), action (reasoning, tool use, and environment interaction), and decision-making (control flow determining when to remember, reason, or act). This decomposition provides a principled vocabulary for describing and comparing the diverse memory architectures surveyed in this section.

\subsubsection{Memory Taxonomy in CoALA}

CoALA distinguishes four memory types, each mapping to established cognitive science constructs and implemented through specific computational mechanisms.

Working memory corresponds to the model's current context window holding active information, mapping to Baddeley's working memory model. In CoALA's formalization, working memory has three key properties: limited capacity (bounded by the context window), rapid access (available to the model's attention mechanism without retrieval), and transience (contents change with each interaction). The central executive function---deciding what information to maintain in working memory---is implemented by the language model itself, which selects what to attend to and what to request from external memory.

Episodic memory stores records of specific interactions with temporal context, corresponding to Tulving's episodic memory. CoALA's episodic memory is implemented through databases, vector stores, or structured logs that preserve the temporal ordering and contextual associations of individual experiences. Retrieval from episodic memory uses similarity-based search (finding memories similar to the current query), temporal proximity (finding memories near a given time), or associative retrieval (following links between related memories).

Semantic memory stores general knowledge and patterns, corresponding to Tulving's semantic memory. In CoALA, semantic memory has two implementations: model parameters (knowledge encoded during pre-training, accessed through the forward pass) and explicit knowledge stores (knowledge graphs, structured databases, or indexed text corpora accessed through retrieval). The dual implementation reflects the CLS framework's distinction between neocortical knowledge (parameters) and explicitly stored knowledge (external stores).

Procedural memory encodes knowledge of how to perform actions, implemented through tool descriptions, workflow definitions, and action templates. This memory type, often overlooked in LLM memory research, is essential for agentic systems that must know not only what but how---how to use tools, follow procedures, and execute multi-step workflows. Procedural memory in CoALA corresponds to production rules in classical cognitive architectures.

\subsubsection{Action Space and Memory Control}

CoALA's action space encompasses internal actions (reasoning steps, memory retrieval, memory storage, summarization) and external actions (tool calls, API invocations, environment interactions). The decision-making component---implemented as a control loop that the language model itself executes---determines when to store information to memory, when to retrieve from memory, when to reason internally, and when to act externally. This illuminates a key design dimension: the degree to which memory management is explicitly controlled by the agent (as in MemGPT, where the LLM issues explicit memory function calls) versus implicitly managed by the architecture (as in attention-based retrieval, where the model's attention mechanism automatically selects relevant information). CoALA suggests that the most effective agents require both explicit memory operations (for deliberate storage and retrieval) and implicit memory operations (for automatic context-based information access), mirroring the distinction between controlled and automatic processing in human cognition.

\subsubsection{Connection to Classical Cognitive Architectures}

CoALA's design draws explicitly on Soar and ACT-R, two classical cognitive architectures that have been developed over decades. Soar organizes cognition around a problem space with operators applied to states, using chunking to compile frequently used reasoning patterns into efficient productions. ACT-R implements a modular architecture with declarative memory (facts), procedural memory (productions), and perceptual-motor modules, with retrieval governed by an activation function combining base-level activation (recency and frequency of use) and spreading activation (contextual relevance). CoALA adapts these concepts for LLM agents: Soar's problem-space formulation maps to the language model's reasoning process, ACT-R's declarative memory maps to CoALA's episodic and semantic stores, and ACT-R's activation-based retrieval maps to the similarity and recency-based retrieval scoring used by systems like Generative Agents. The key innovation of CoALA is recognizing that an LLM can serve as the cognitive controller that classical architectures implemented through hand-crafted rules.

\subsection{Generative Agents}
\label{subsec:generative_agents}

Park et al.\ \cite{park2023generative} demonstrated human-inspired memory architecture through 25 autonomous agents in a simulated town environment. Each agent maintains a memory stream---a comprehensive record of observations, conversations, and actions, timestamped and stored as natural language---serving as episodic memory.

\subsubsection{Memory Stream Architecture}

The memory stream stores observations with metadata including timestamp, location, participants, and emotional valence. Retrieval is governed by a scoring function combining recency (exponential decay), importance (significance rated 1--10 by the LLM), and relevance (semantic similarity via embedding cosine similarity). This three-factor function implements principles from cognitive science: the recency effect, importance-based encoding from emotional memory theory, and associative retrieval from spreading activation models.

\subsubsection{Reflection and Planning}

Reflection is a periodic process where the agent reviews recent memories and synthesizes higher-level insights stored as new memories, implementing memory consolidation from specific episodes to generalizable semantic knowledge. Planning uses memories and reflections to generate structured plans decomposed into hourly and sub-hourly actions, creating a feedback loop between memory, reflection, and action that mirrors the human process of experience-guided behavior.

\subsubsection{Emergent Social Behavior}

The most striking finding is that this memory architecture---timestamped observation log with recency-importance-relevance retrieval, periodic reflection, and plan generation---produces sophisticated emergent social behaviors: forming relationships, spreading information through social networks, coordinating collective activities, and referencing past interactions contextually. Agents spontaneously organized a Valentine's Day party through social coordination that no single agent planned: one agent mentioned the idea, others spread it through conversations, and the group collectively organized decorations, invitations, and timing. This emergence suggests that human memory organization principles are functional requirements for social intelligence, not merely metaphorical parallels.

\subsubsection{Architectural Significance}

The Generative Agents architecture is significant beyond its specific application to simulated environments because it demonstrates several general principles. First, relatively simple memory mechanisms (timestamped storage with three-factor retrieval scoring) can produce sophisticated behavior when operating within a well-designed cognitive loop. Second, the combination of bottom-up episodic storage (recording observations as they occur) with top-down reflective processing (periodically synthesizing higher-level insights) creates a bidirectional information flow that enriches both the episodic and semantic levels. Third, the architecture's success with a 4K context window demonstrates that intelligent memory management can substitute for larger context windows: the system processes information far exceeding 4K tokens by loading only the most relevant memories for each decision point. This principle of intelligent loading over brute-force inclusion is central to the argument that hierarchical memory architectures are necessary for scalable context management.

\subsection{MemoryBank and Forgetting-Curve Memory Management}

MemoryBank \cite{zhong2024memorybank} implements explicit Ebbinghaus forgetting curve dynamics. Each stored memory has a strength decaying exponentially, refreshed upon retrieval (mimicking the spacing effect). Memories below a strength threshold are archived (compressed) or deleted, implementing adaptive forgetting.

\subsubsection{Dual Storage Architecture}

MemoryBank maintains short-term memory (recent interactions at full fidelity) and long-term memory (older interactions summarized and compressed). The transition is governed by the forgetting curve: aging memories are candidates for summarization, while frequently-retrieved high-strength memories resist compression. This directly implements the cognitive distinction between detailed episodic memories and schematic semantic memories.

\subsubsection{SiliconFriend Application}

MemoryBank was demonstrated through SiliconFriend, a long-term companion chatbot maintaining persistent memory across sessions, tracking preferences, remembering shared experiences, and adapting behavior. Evaluations showed forgetting-curve-based management produced more natural responses than sliding-window or complete-logging baselines.

\subsection{Self-Controlled Memory Framework}

The Self-Controlled Memory (SCM) framework \cite{liang2023scm} introduces dynamic memory management with explicit short-term and long-term storage managed by the LLM as memory controller. SCM processes ultra-long input streams iteratively, reading manageable chunks and using the LLM to determine what to retain.

\subsubsection{Flash and Archived Memory}

Flash memory holds recent, important information with limited capacity (approximately 1K tokens). Archived memory holds older information in summarized form with unlimited capacity. The LLM-based memory controller evaluates each chunk against current memory state, storing important information in flash memory and compressing older contents to archived memory.

\subsubsection{Performance}

SCM achieves 78\% retrieval recall versus 45\% for baselines on long-term dialogue tasks and handles texts of 100K or more tokens despite a 4K--8K context window, demonstrating that LLM-directed memory management can effectively simulate much larger context windows.

\subsection{Hierarchical Memory Transformer}

The Hierarchical Memory Transformer (HMT), published at NAACL~2025, implements a three-tier memory hierarchy within the transformer architecture, inspired by the Atkinson-Shiffrin model.

\subsubsection{Three-Tier Architecture}

Sensory memory stores key-value embeddings from early input in a fixed-capacity buffer, analogous to the sensory register. Working memory maintains full self-attention within a sliding window (512--2K tokens), corresponding to short-term memory with limited-capacity but high-fidelity processing. Current focus computes attention over both working memory and selected sensory memory entries using sparse patterns, integrating local and global information.

\subsubsection{Efficiency and Performance}

HMT achieves $O(w^2 + n \cdot k)$ complexity providing 2.5--116$\times$ speedup over full attention on 32K-token contexts while maintaining comparable short-context quality and superior long-context quality. The architecture demonstrates that the Atkinson-Shiffrin hierarchy provides an efficient computational architecture.

\subsection{Strata: Hierarchical Context Caching}

Strata applies hierarchical memory principles to KV cache management in production serving: GPU VRAM (hot cache, 1K--4K tokens, full precision, immediate access), CPU RAM (warm cache, 4K--32K tokens, compressed FP8), and disk/NVMe (cold cache, 32K--1M+ tokens, highly compressed). Cache management follows LRU eviction with promotion and demotion between tiers, mirroring the biological memory hierarchy where the GPU hot cache corresponds to working memory, CPU warm cache to short-term memory, and disk cold cache to long-term memory.

\subsection{RecallM, RET-LLM, and TransformerFAM}

RecallM \cite{kynoch2023recallm} introduces temporal understanding in memory, considering temporal relationships (before, after, during) when scoring relevance, enabling answers about event ordering and duration. RET-LLM \cite{modarressi2024retllm} proposes general-purpose read-write memory learning to perform memory operations through training on memory-intensive tasks. TransformerFAM \cite{irie2024transformerfam} proposes that feedback attention---output fed back as input to the same layer---functions as working memory analogous to Baddeley's phonological loop.

\subsection{RAISE and Cognitive LLMs}

RAISE \cite{liu2024raise} presents a memory-enhanced architecture with dual memory (short-term conversational context, long-term user model) connected by consolidation operations. Recent work on Cognitive LLMs integrates classical cognitive architectures (ACT-R-like production systems) with LLMs for manufacturing decision-making, using symbolic production rules and working memory buffers alongside LLM flexibility. Decision-making processes are encoded as latent representations injected via LoRA \cite{hu2022lora}, achieving benefits in interpretability, knowledge transfer, reliability, and explainability.

\subsubsection{Global Workspace Theory}

The Cognitive Workspace hypothesis applies Baars's Global Workspace Theory to LLM design. GWT posits conscious processing occurs in a limited-capacity workspace where cognitive processes compete for access. The LLM implementation maintains a workspace (1K--4K tokens) containing current goals, observations, retrieved knowledge, and reasoning traces. Multiple cognitive modules (reading, planning, memory, execution) operate in parallel with the workspace as shared communication channel, providing principled guidelines for context window allocation.

\subsubsection{Multiple Memory Systems for Agents}

Inspired by cognitive psychology, recent work implements four distinct memory types for agents: episodic (events with timestamps and context), semantic (structured facts in knowledge graphs), procedural (learned skills and action sequences), and implicit (habits and preferences influencing behavior without explicit retrieval). This multi-system approach enables specialized encoding, storage, and retrieval for each knowledge type, aligning with the functional requirements of different information categories.

\subsection{Comparative Analysis: Human versus LLM Memory}

\subsubsection{Shared Phenomena}

Both humans and LLMs exhibit serial position effects \cite{murdock1962serial, glanzer1966two}: primacy effects, recency effects, and the U-shaped serial position curve. Both produce false memories or hallucinations---humans through the Deese-Roediger-McDermott paradigm, LLMs through semantic contamination. Both exhibit capacity-limited working memory: humans maintain approximately 4--7 chunks, while LLMs show empirically similar limitations in simultaneous reasoning capacity.

\subsubsection{Critical Differences}

Human memory is highly selective through amygdala-hippocampal emotional tagging; LLM memory is non-selective, treating all tokens equally. Human forgetting is adaptive, with rates sensitive to importance; LLM forgetting is non-adaptive, occurring through truncation or catastrophic forgetting. Human concepts are stable across contexts; LLM concept representations are context-dependent. Most critically, human metacognition is well-calibrated (JoL-recall correlation $r \approx 0.68$), while LLM metacognition is poorly calibrated ($r = 0.12$--$0.18$ across major models), meaning LLMs cannot reliably assess their own output accuracy.

\subsubsection{Design Implications}

These differences suggest design principles: importance tagging for selective encoding, adaptive forgetting replacing arbitrary truncation, metacognitive calibration for reliability assessment, and explicit episodic structures with timestamps and associative retrieval cues. These principles converge on the conclusion that human-inspired hierarchical memory is a functionally necessary architecture.

\subsection{Catastrophic Forgetting and Continual Learning}
\label{subsec:catastrophic}

Catastrophic forgetting---the phenomenon where training on new data causes severe degradation on previously learned tasks---represents the computational analogue of retroactive interference in human memory, where new learning impairs recall of previously learned material. However, the severity differs dramatically: while humans experience modest retroactive interference that is largely offset by consolidation mechanisms, neural networks experience catastrophic interference that can completely overwrite previously learned knowledge.

\subsubsection{Empirical Characterization}

Luo et al.\ \cite{luo2023empirical} provided a systematic empirical characterization of catastrophic forgetting in large language models, documenting that sequential fine-tuning on multiple tasks causes significant performance degradation on earlier tasks. The degradation follows predictable patterns: tasks requiring factual recall are most vulnerable (with up to 80\% accuracy drops), tasks requiring reasoning are moderately affected, and tasks requiring linguistic competence are least affected (suggesting that syntactic knowledge is more deeply encoded than factual knowledge). This pattern mirrors findings in human neuropsychology, where semantic knowledge (deeply encoded through repeated exposure) is more resistant to interference than episodic memories (encoded through single exposures).

\subsubsection{Parameter-Efficient Mitigation}

LoRA \cite{hu2022lora} addresses catastrophic forgetting by restricting fine-tuning updates to low-rank subspaces of the model's weight matrices: $W_{\text{new}} = W_{\text{base}} + BA$ where $B \in \mathbb{R}^{d \times r}$ and $A \in \mathbb{R}^{r \times d}$ with $r \ll d$. This implements a functional separation between semantic memory (frozen base parameters $W_{\text{base}}$, encoding general knowledge from pre-training) and task-specific adaptations (low-rank updates $BA$, encoding new task knowledge). The separation directly parallels the complementary learning systems framework: the base parameters function as the slowly-learned neocortical system, while LoRA adapters function as quickly-learned hippocampal representations that do not interfere with the base knowledge. Multiple LoRA adapters can be trained for different tasks and swapped at inference time, providing a form of procedural memory switching.

\subsubsection{Continual Learning Strategies}

Broader continual learning strategies implement increasingly sophisticated mechanisms for balancing new learning with knowledge preservation. Elastic Weight Consolidation (EWC) penalizes changes to parameters identified as important for previous tasks (using the Fisher information matrix as an importance measure), implementing a form of synaptic consolidation where well-established neural connections resist modification. Experience replay intermixes training data from old and new tasks, preventing the distributional shift that causes forgetting---this directly analogizes hippocampal replay during sleep, where memories of past experiences are replayed to maintain their neural representations. Progressive neural networks add new capacity (columns) for each new task while freezing parameters for old tasks, implementing a growth-based strategy that avoids interference entirely at the cost of increasing model size. PackNet prunes unnecessary connections after each task, freeing parameters for future tasks while preserving essential connections.

These strategies provide computational mechanisms for the memory consolidation and interference prevention observed in biological systems, and their success validates the cognitive science insight that effective long-term learning requires explicit mechanisms to protect existing knowledge from disruption by new experience.

\subsection{Toward Unified Multi-Scale Memory}

The surveyed architectures converge on a unified principle: effective context management requires memory organized across multiple temporal scales and abstraction levels, with mechanisms for transferring information between levels. This convergence is remarkable because the architectures were developed by independent research groups solving different problems (simulation, conversation, document processing, agent design) yet arrived at structurally similar solutions, suggesting that the hierarchical principle is a fundamental design constraint rather than an arbitrary choice.

\subsubsection{The Five-Level Memory Hierarchy}

The unified framework that emerges from the surveyed systems comprises five levels, each characterized by temporal scale, capacity, access speed, and representational fidelity.

The immediate level (microseconds, 1--2 positions) corresponds to the current attention focus---the specific tokens receiving the highest attention weights at the current generation step. This level corresponds to the sensory register in the Atkinson-Shiffrin model and the current focus of attention in Baddeley's model. In computational terms, it is implemented by the softmax attention output at the current position.

The working level (seconds to minutes, 512--4K tokens) corresponds to the local context maintained with full attention. This level holds the information actively being processed, including the current query, recent dialogue turns, and any retrieved information loaded for the current task. It corresponds to Baddeley's working memory and is implemented by the model's context window contents.

The short-term level (minutes to hours, 4K--32K tokens) holds recent events in extended context or KV cache. This level contains information from the current session that has not yet been compressed or archived, implemented through extended context windows (Section~\ref{sec:window}), sliding window attention, or warm KV cache (as in Strata).

The long-term episodic level (hours to days) stores compressed and summarized representations in external memory, accessible through retrieval. This corresponds to Tulving's episodic memory and is implemented through vector databases, structured logs, or summarized conversation history.

The semantic level (permanent) encodes general knowledge in model parameters, corresponding to Tulving's semantic memory and McClelland's neocortical system. This knowledge is accessed through the model's forward pass without explicit retrieval.

This framework moves from faster access and greater specificity at the top to larger capacity and greater generality at the bottom, mirroring both the Atkinson-Shiffrin hierarchy and biological memory. Information flows upward through retrieval (loading relevant long-term memories into working memory) and downward through consolidation (compressing working memory contents into long-term storage).

\subsubsection{Inter-Level Transfer Mechanisms}

The effectiveness of the hierarchical framework depends critically on mechanisms for transferring information between levels. Upward transfer (retrieval) brings information from lower, larger-capacity levels into the constrained working memory where it can be processed by the language model. The retrieval mechanism determines what information is available for processing and therefore constrains the system's capabilities. Effective retrieval requires relevance scoring (finding information pertinent to the current task), temporal awareness (understanding when stored information was created and whether it remains current), and importance weighting (prioritizing high-value information when working memory capacity is limited).

Downward transfer (consolidation) moves information from higher, more detailed levels into lower, more compressed levels for long-term preservation. Consolidation mechanisms include summarization (compressing detailed episodes into schematic summaries), abstraction (extracting general principles from specific experiences), and indexing (creating retrieval cues that enable future access to consolidated memories). The timing of consolidation is critical: too-frequent consolidation loses detail prematurely, while too-infrequent consolidation allows working memory to overflow.

\subsubsection{Mapping Existing Systems to the Hierarchy}

Each surveyed system implements a subset of this hierarchy. MemGPT implements working memory (context window), short-term memory (working context buffer), and long-term memory (external databases) with explicit LLM-controlled transfer. Generative Agents implement working memory (current context), episodic memory (memory stream), and a consolidation mechanism (reflection). MemoryBank implements short-term memory (recent interactions) and long-term memory (summarized history) with Ebbinghaus-curve-based transfer. The Hierarchical Memory Transformer implements sensory memory (KV buffer), working memory (sliding window attention), and selective retrieval from the sensory buffer. Strata implements the hierarchy at the hardware level: GPU VRAM (hot/working), CPU RAM (warm/short-term), and disk (cold/long-term). No single existing system implements all five levels with principled inter-level transfer, representing a clear direction for future research.


% ============================================================
% SECTION 8: HALLUCINATION AND CONTEXT FAILURES
% ============================================================
\section{Hallucination as Context Management Failure}
\label{sec:hallucination}

Hallucination---generation of fluent but factually incorrect or unsupported text---represents a significant practical challenge for deployed language models. While hallucination has multiple causes including training data contamination, knowledge gaps, and decoding artifacts, a substantial class arises from context management failures: the model fails to attend to, retrieve, or correctly integrate needed information from the provided context. This section examines the taxonomy of hallucinations, the specific context management mechanisms that produce them, the interaction between retrieval-augmented generation and hallucination, and detection and mitigation strategies.

\subsection{Taxonomy of Hallucinations}

The foundational taxonomy of hallucinations was established by Maynez et al.\ \cite{maynez2020faithfulness}, who conducted a large-scale human evaluation of neural abstractive summarization systems on the XSum dataset. Their framework distinguishes two fundamental categories based on the relationship between the hallucinated content and the source material.

\subsubsection{Intrinsic Hallucination}

Intrinsic hallucination occurs when the model's output contradicts the source document---the generated content is provably false given the provided context. This represents a direct context utilization failure: the model has access to the correct information but generates contradictory output. Intrinsic hallucinations can be detected through natural language inference (NLI) entailment checking, where a contradiction between the source and the output indicates hallucination. Maynez et al.\ found that intrinsic hallucinations are frequent across all abstractive summarization models and correlate with model quality: pre-trained models (BART, PEGASUS) produce fewer intrinsic hallucinations than standard sequence-to-sequence models, and textual entailment-based metrics correlate better with faithfulness than standard overlap metrics like ROUGE.

\subsubsection{Extrinsic Hallucination}

Extrinsic hallucination occurs when the model's output contains information absent from the source document---content that cannot be verified from the provided context. Critically, extrinsic hallucinations are not necessarily factually incorrect: the model may generate externally accurate information that is simply not grounded in the source. For example, a summary might correctly state a public figure's educational background even though the source article does not mention it. This nuance has important implications for context management: extrinsic hallucinations represent context sufficiency failures (the model draws on parametric knowledge rather than the provided context) rather than context utilization failures (the model misuses provided information).

Ji et al.\ \cite{ji2023survey} and Huang et al.\ \cite{huang2023survey} extended this taxonomy to broader language generation tasks, identifying additional hallucination types including factual fabrication (generating plausible but false claims), entity confusion (conflating attributes of similar entities), and temporal confusion (mixing information from different time periods). These extensions are directly relevant to context management because each failure mode can be traced to specific attention or retrieval failures.

\subsection{Context-Dependent Hallucination Mechanisms}

Several specific context management mechanisms systematically produce hallucinations. Understanding these mechanisms is essential for designing systems that minimize hallucination risk.

\subsubsection{Lost-in-the-Middle Hallucinations}

The lost-in-the-middle phenomenon \cite{liu2024lost} directly causes hallucination when needed information is positioned in the middle of context and receives insufficient attention. When the model cannot access relevant information due to positional attention deficits, it may fabricate plausible answers from parametric knowledge, misattribute information from context boundaries to middle-positioned queries, or generate responses that are locally coherent but factually unsupported by the provided context. This failure mode is particularly insidious in RAG systems, where the ordering of retrieved documents within the prompt directly determines whether the model attends to the evidence supporting the correct answer. Research has shown that placing the most relevant document in the middle of 20 retrieved documents can reduce answer accuracy by over 50\% compared to placing it at the beginning or end.

\subsubsection{Positional Biases and Serial Position Effects}

Primacy bias influences approximately 10\% of decisions across models, amplified for multiple-choice questions and only partially mitigated by model scaling. The bias operates through the attention sink mechanism described in Section~\ref{sec:attention}: initial tokens absorb disproportionate attention, causing the model to over-weight information presented early in the context. Recency bias similarly affects ranking tasks where recently presented information receives disproportionate weight, driven by the causal attention mask that gives later tokens fewer context elements to attend to and thus more concentrated attention on recent predecessors.

The combined effect creates a U-shaped vulnerability curve: reliable attention at context boundaries, systematic underutilization in the middle. For contexts exceeding 10K tokens, middle content becomes largely inaccessible to the model's attention mechanism. These biases have been confirmed across all major model families including GPT-4, Llama, and Mistral variants. Emerging research has also established that LLMs can induce cognitive biases in users through positional effects in generated text, creating feedback loops where biased generation leads to biased interpretation.

\subsubsection{Context Conflict and Contradiction}

When the provided context contains contradictory information---whether from multiple retrieved documents that disagree, from context that contradicts the model's parametric knowledge, or from ambiguous phrasing---the model must resolve the conflict to generate a coherent response. Research has shown that models tend to favor their parametric knowledge over provided context when conflicts arise, particularly when the parametric knowledge is well-established. This preference hierarchy can cause hallucination when the provided context is correct but contradicts the model's (incorrect) parametric beliefs. The interaction between context and parametric knowledge remains poorly understood and represents an active research area.

\subsection{Hallucination in Retrieval-Augmented Generation}

RAG systems were developed partly to mitigate hallucination by grounding generation in retrieved evidence, but they introduce their own characteristic failure modes that can produce novel hallucination patterns.

\subsubsection{Retrieval-Phase Hallucination Sources}

Hallucinations in RAG can originate during the retrieval phase through several mechanisms. Poor retrieval quality---retrieving irrelevant, outdated, or misleading passages---provides the model with incorrect or unhelpful context, potentially causing it to generate answers grounded in irrelevant evidence. Incomplete retrieval misses relevant passages, leaving the model without the information needed to answer correctly and forcing it to rely on parametric knowledge. Query-document mismatch occurs when the retrieval query does not accurately capture the information need, particularly for complex multi-hop queries that require synthesizing information across multiple passages.

\subsubsection{Generation-Phase Hallucination Sources}

Even with high-quality retrieval, hallucinations can arise during generation. Context noise occurs when retrieved passages contain distracting information alongside relevant content, causing the model to incorporate irrelevant details. The ``middle curse'' (the lost-in-the-middle effect applied to retrieved documents) means that evidence positioned in the middle of the retrieved document set may be underutilized. Alignment problems arise when the generation does not match the query intent despite relevant retrieved context, often caused by the model following its pre-training distribution rather than the specific evidence provided.

\subsubsection{Domain-Specific Hallucination Risks}

Research on legal RAG systems has demonstrated that hallucinations persist even with high-quality retrieval in specialized domains, with document complexity affecting hallucination rates and citation accuracy being particularly challenging. In healthcare applications, RAG systems face elevated hallucination risks due to the precision required in medical contexts, where even minor factual errors can have serious consequences. These domain-specific findings underscore that generic RAG architectures require domain-specific calibration and verification layers to achieve acceptable hallucination rates in high-stakes applications.

\subsection{Hallucination Detection and Mitigation}

\subsubsection{Detection Methods}

TruthfulQA \cite{lin2022truthfulqa} evaluates truthfulness against common misconceptions, testing whether models generate factually correct responses rather than plausible-sounding but incorrect ones. HaluEval \cite{li2023halueval} provides large-scale hallucination detection evaluation across QA, dialogue, and summarization tasks. FActScore \cite{min2023factscore} enables fine-grained factual precision evaluation by decomposing text into atomic claims and verifying each against reference knowledge, achieving reliable automated evaluation that correlates well with human judgment.

\subsubsection{RAG-Based Mitigation}

RAG systems mitigate hallucination by providing source material that grounds generation in evidence, but the effectiveness depends on retrieval quality. Advanced RAG variants address specific hallucination failure modes. Self-RAG \cite{asai2024selfrag} reduces hallucination by having the model decide whether retrieval is needed and evaluating whether generated content is supported by retrieved evidence. CRAG \cite{yan2024corrective} addresses poor retrieval quality through confidence-based corrective actions. GraphRAG provides structured multi-hop reasoning that reduces hallucination on complex queries requiring information synthesis.

\subsubsection{Generation-Level Mitigation}

Citation generation \cite{gao2023enabling} constrains models to generate only claims explicitly grounded in cited evidence, providing traceability from claims to sources. Self-consistency \cite{wang2023selfconsistency} samples multiple reasoning chains and selects the most frequent answer, exploiting the insight that hallucinations are locally plausible but globally inconsistent---different reasoning paths are unlikely to converge on the same hallucinated answer. Chain-of-Verification \cite{dhuliawala2023chain} generates verification questions, answers them independently, and revises the original response based on detected inconsistencies. CREAM-RAG introduces improved passage filtering, better context selection, and enhanced generation constraints specifically designed to limit hallucination in RAG pipelines.

\subsubsection{Groundedness in Long-Form Generation}

Research on groundedness in retrieval-augmented long-form generation has revealed that even with retrieval, hallucinations persist in long-form outputs, with the beginning of generated text being more grounded in retrieved evidence than later portions. This temporal degradation in groundedness suggests that the model's reliance on retrieved context diminishes as generation proceeds, with later tokens increasingly drawing on parametric knowledge rather than provided evidence. Explicit grounding mechanisms---constraining each generated sentence to be traceable to a specific retrieved passage---improve groundedness but at the cost of generation fluency and coherence.

\subsubsection{The FRAMES Evaluation Framework}

The FRAMES (Fact, Fetch, and Reason) benchmark provides a unified evaluation of RAG systems across three dimensions: factuality (accuracy of final answers), retrieval quality (relevance and coverage of retrieved passages), and reasoning (correctness of multi-step inference). This multi-dimensional evaluation enables identification of specific failure points in RAG pipelines, distinguishing hallucinations caused by poor retrieval from those caused by poor reasoning over correctly retrieved evidence.


% ============================================================
% SECTION 9: MULTI-TURN DIALOGUE
% ============================================================
\section{Multi-Turn Dialogue and Conversational Memory}
\label{sec:multiturn}

Multi-turn dialogue presents unique context management challenges beyond single-query processing. Conversations evolve over multiple turns, sessions, and potentially months, requiring coherent understanding of ongoing dialogue, entity state tracking, reference resolution, and cross-session continuity. Unlike document processing where the entire input is available from the start, dialogue is inherently incremental: information arrives turn by turn, must be integrated with existing understanding, and may require revision as the conversation progresses. This section examines the specific challenges of multi-turn context management, the architectures developed to address them, dialogue state tracking as a structured approach to conversational memory, and the evaluation methodologies for assessing long-term conversational coherence.

\subsection{Challenges in Multi-Turn Context Management}

\subsubsection{Context Growth and Window Exhaustion}

As conversations progress, accumulated context grows continuously---a typical dialogue generating 100--500 tokens per turn may consume 10,000--50,000 tokens after 100 turns, and conversations extending over multiple sessions can generate hundreds of thousands of tokens. For customer service interactions, technical support sessions, or therapeutic conversations, the relevant history may span weeks or months. This progressive context growth inevitably exceeds the model's context window, requiring either truncation (losing early context), compression (losing detail), or external storage (adding retrieval latency). Each approach introduces characteristic failure modes: truncation causes the system to ``forget'' early agreements or commitments; compression may distort nuances in emotional tone or implicit commitments; retrieval may miss contextually relevant but semantically distant memories.

\subsubsection{Coreference and Ellipsis Resolution}

Pervasive use of pronouns, ellipsis, and references to earlier turns requires fine-grained access to specific utterances. A statement like ``Can you change that to Tuesday instead?'' requires resolving ``that'' to a specific earlier commitment, ``Tuesday'' to the correct week, and ``instead'' to the original arrangement---potentially requiring access to turns from minutes or hours earlier. Research has demonstrated that approximately 79\% of natural conversations include explicit references to earlier episodes, underscoring that fine-grained episodic memory access is not an edge case but a core requirement for natural dialogue.

\subsubsection{Dynamic Entity State Tracking}

Entity states change over dialogue: orders progress from ``pending'' to ``shipped'' to ``delivered''; user preferences evolve as they learn more about options; commitments are made, modified, and sometimes retracted. Tracking these state changes requires maintaining structured representations that are updated as new information arrives, creating a dynamic knowledge base that reflects the current state of all relevant entities. This structured tracking goes beyond simple memory retrieval: it requires interpreting new utterances in the context of existing state and applying appropriate updates.

\subsubsection{Cross-Session Continuity}

For persistent assistants, conversations span sessions separated by days or months, requiring persistent storage with activation mechanisms for new sessions. A user returning after a week expects the system to remember prior discussions, agreed-upon preferences, and outstanding tasks without explicit prompting. This requires not only storing past interactions but recognizing when stored information is relevant to a new session---a form of proactive memory retrieval that anticipates information needs rather than responding to explicit queries.

\subsection{Conversation Memory Architectures}

Multiple architectural approaches address these challenges, each representing different trade-offs between memory fidelity, computational cost, and retrieval effectiveness.

\subsubsection{Buffer-Based Approaches}

The simplest approach maintains recent turns as a fixed-size window, typically the last $N$ messages (where $N$ ranges from 5 to 20 in practical systems). This message windowing approach works well for short conversations where all relevant context is recent, but performance degrades as important context from earlier messages is lost. The computational cost is minimal, making it suitable for high-throughput applications where context retention is secondary to response latency. Growing buffer approaches expand the context window as the conversation progresses, eventually hitting the model's context limit and requiring fallback to other strategies.

\subsubsection{Summary-Based Memory}

Summary-based memory uses the LLM to generate running summaries of conversation history, trading detail for capacity. Rolling summarization triggers after a configurable number of turns (typically 5--10), with the LLM creating a concise summary that replaces the original messages. Subsequent summaries incorporate the prior summary alongside new context, creating a recursive chain of increasingly compressed history. Research on recursive summarization for long-term dialogue memory, published at ICLR~2025, demonstrates that this recursive approach enables at least 5$\times$ more information to be retained in the chat context compared to simple truncation. However, rolling summarization suffers from cascading information loss: each summarization step discards detail, and errors in early summaries propagate to subsequent ones, with nuances and subtle shifts in intent being lost over time.

\subsubsection{Hierarchical Summary Approaches}

Hierarchical summary approaches maintain multiple granularities: recent turns at full fidelity, older turns as per-turn summaries, and still older turns as session-level summaries, mirroring the human memory hierarchy's transition from detailed short-term to schematic long-term representations. H-MEM implements a four-layer memory structure (Domain, Category, Memory Trace, Episode) with position-indexed retrieval from each level. SGMEM organizes conversation memory as a sentence-level graph with explicit semantic associations, enabling structured retrieval and addressing memory fragmentation that occurs with simple list-based storage. These hierarchical approaches achieve compression ratios of up to 10$\times$ while maintaining structured retrieval capability.

\subsubsection{Vectorized Memory}

Vectorized memory stores interactions as vector embeddings in a database, retrieving semantically similar past conversations when relevant. This approach excels at finding topically related past exchanges but may miss important context that is semantically dissimilar to the current query (such as a commitment made during an unrelated discussion). Combining vectorized retrieval with temporal proximity and importance weighting---as in the Generative Agents retrieval paradigm \cite{park2023generative}---provides more robust recall:

\begin{equation}
\text{Score}(m) = \alpha \cdot \text{Recency}(m) + \beta \cdot \text{Relevance}(m, q) + \gamma \cdot \text{Importance}(m)
\end{equation}

where $\alpha, \beta, \gamma$ are tunable weights balancing temporal proximity, semantic relevance to the current query $q$, and rated importance of the memory $m$.

\subsubsection{Proactive Memory Extraction}

Recent work goes beyond static summarization to proactive memory extraction, where the system anticipates which information from current interactions will be useful for future tasks and proactively stores it in structured memory. This approach, demonstrated in agent systems, enables the memory module to identify and extract action-relevant information (user preferences, commitments, task progress) rather than passively storing verbatim history or generic summaries. Proactive extraction represents a shift from reactive memory (storing everything and retrieving on demand) to goal-directed memory (strategically storing information likely to be needed).

\subsection{Dialogue State Tracking}

Dialogue State Tracking (DST) provides a structured approach to conversational memory in task-oriented dialogue systems, maintaining a formal representation of the conversation state as a set of slot-value pairs that capture user goals, gathered information, and remaining tasks.

\subsubsection{Core Architecture}

A DST system maintains predefined slots for each domain (for a restaurant booking system: cuisine, location, price range, number of guests, date, time) and updates slot values based on each user utterance. The state tracker interprets user input in the context of the existing state, resolving ambiguous references (``same place as last time'') and applying corrections (``actually, make it Italian instead''). The tracked state drives dialogue policy decisions (which information to request next) and external actions (API calls, database queries).

\subsubsection{Multi-Domain DST}

Real-world conversations often span multiple domains (booking a flight, then a hotel, then a restaurant), requiring the state tracker to handle domain transitions, resolve conflicting slot definitions across domains, and determine which information carries forward between domains. Knowledge graph-based approaches frame multi-domain DST as graph question answering, naturally handling cross-domain entity relationships. Modern approaches using large language models as state updaters achieve better zero-shot performance on unseen domains, reduced need for extensive training data, and more natural handling of out-of-vocabulary inputs, though they face challenges with hallucination in slot value extraction and maintaining consistency across long conversation histories.

\subsubsection{Chain-of-Thought DST}

Chain-of-Thought approaches incorporate explicit reasoning chains in state tracking, improving both accuracy and interpretability. Rather than directly predicting slot values from the dialogue, the model generates a reasoning trace explaining why each slot update is warranted, enabling verification of the tracking logic and identification of errors. This approach achieves improved accuracy on challenging multi-turn scenarios where implicit context must be resolved through reasoning.

\subsection{User Modeling and Personalization}

Long-term conversational agents must not only remember what was discussed but learn about the user as a person---their preferences, interests, communication style, expertise level, and goals. User modeling represents a distinct memory challenge that goes beyond episodic recall to construct and maintain a persistent user profile that evolves over time.

\subsubsection{Dynamic User Profiling}

Research on dynamic user profiling reveals that frontier models including GPT-4, DeepSeek-R1, Gemini-2.0, and Llama-4 struggle significantly with user awareness---the ability to apply knowledge about a user's preferences and traits to new scenarios. The PersonaMem benchmark, containing over 180 simulated user profiles with up to 60 sessions per user across 15 real-world tasks, demonstrates a significant performance gap when knowledge about user preferences must be applied across novel scenarios. Models can often recall stated preferences when directly asked but fail to proactively apply those preferences when generating responses to new requests, suggesting that user modeling requires not just storage but active integration of user knowledge into response generation.

\subsubsection{Preference Evolution}

User preferences are not static: they evolve based on experience, context, and time. Research on teaching language models to evolve with users demonstrates that dynamic profile modeling---tracking how preferences change over time rather than maintaining a fixed profile---significantly improves personalization quality. This temporal dimension of user modeling connects to the broader memory management challenge: the system must distinguish between stable traits (communication style, expertise level) and evolving preferences (favorite tools, current projects, shifting interests), applying different retention and update strategies to each.

\subsubsection{Curiosity-Driven Personalization}

An emerging approach incorporates curiosity-based intrinsic reward in multi-turn training, encouraging the model to actively learn about users through strategic questioning and observation. Rather than passively recording stated preferences, the curiosity-driven agent optimizes for improving its user model accuracy over time, balancing immediate response quality with long-term learning about the user. This active learning approach treats user modeling as an ongoing inference problem rather than a simple storage task, aligning with cognitive science research on theory of mind development.

\subsection{Context Window Budgeting}

In production systems, the context window must be shared among multiple competing demands: system prompts, conversation history, retrieved documents, tool outputs, and space reserved for generation. Context budgeting---strategic allocation of available tokens across these components---is critical for maximizing information utilization within fixed window constraints.

\subsubsection{Token Allocation Strategies}

Practical context budgeting divides the available window into regions with priority-based allocation. System prompts (defining agent behavior, persona, and constraints) typically consume 500--2,000 tokens and are allocated first as non-negotiable. Retrieved documents (RAG context) receive a variable allocation based on query complexity, typically 2,000--8,000 tokens. Conversation history receives the remaining budget, with recent turns at full fidelity and older turns summarized or omitted. Generation space (tokens reserved for the model's output) must be budgeted to prevent truncation of responses. The specific allocation depends on the application: customer service systems prioritize conversation history, research assistants prioritize retrieved documents, and coding assistants prioritize code context.

\subsubsection{Token-Budget-Aware Reasoning}

Recent research on token-budget-aware LLM reasoning demonstrates that models with context awareness---those that track their remaining context window during conversation---enable more effective task execution through dynamic allocation. These models adjust their behavior based on available context: using more detailed reasoning when budget is ample and more concise processing when budget is constrained. This adaptive behavior mirrors human cognitive resource allocation, where deeper processing is reserved for important or difficult tasks while routine tasks receive shallower processing.

\subsubsection{The Context Window Problem in Agent Systems}

For agent systems that perform multi-step reasoning, tool use, and iterative refinement, context exhaustion is a critical failure mode. Each tool call and its output consumes context budget, and complex tasks requiring many steps may exhaust the window before completing. Strategies for managing this include progressive summarization of intermediate results, selective retention of tool outputs (keeping only the relevant portions), and task decomposition (breaking complex tasks into sub-tasks that each fit within the context budget). The context window problem for agents represents one of the most pressing practical challenges in deployed LLM systems, as it directly limits the complexity of tasks that agents can perform.

\subsection{Psychology-Informed Conversational Memory}

Recent work applies psychological models of memory importance and forgetting to chatbot memory, moving beyond engineering heuristics toward principled, cognitively-grounded memory management. Interaction importance is modeled through multiple psychological dimensions: emotional salience (interactions with strong emotional content are remembered better, paralleling amygdala-mediated encoding in biological memory), goal relevance (interactions related to active goals receive preferential encoding), novelty (new or surprising information is encoded more strongly than routine exchanges), and encoding depth (interactions that require deep processing create stronger memories than superficial exchanges). Decay mechanisms following Ebbinghaus's forgetting curve allow low-importance interactions to fade, with retrieval events resetting the decay clock (implementing the spacing effect). This psychologically-informed approach produces more natural conversational behavior than uniform retention or simple recency-based approaches, as it mirrors the selective, importance-weighted nature of human conversational memory.

The integration of psychological memory models into conversational systems represents the convergence of two themes central to this survey: the practical need for efficient memory management in constrained systems and the theoretical insight that human memory organization reflects fundamental principles of effective information management. The success of psychology-informed approaches in producing more natural, coherent, and user-satisfying conversations provides empirical validation of the cognitive science foundations discussed in Section~\ref{sec:hierarchical}.

\subsection{Long-Term Conversational Memory Evaluation}

\subsubsection{The LoCoMo Benchmark}

The LoCoMo benchmark evaluates very long-term conversational memory through multi-session dialogues with specific quantitative characteristics: 300+ turns per dialogue, 35+ sessions, spanning 6--12 simulated months with up to 25 causally-related events and approximately 9,000 tokens per dialogue. LoCoMo evaluates five reasoning types: single-hop retrieval (direct fact recall), multi-hop reasoning (inference across multiple facts), temporal reasoning (time-aware queries), commonsense reasoning (background knowledge integration), and adversarial queries (robustness testing).

Results reveal the current state of conversational memory capabilities: long-context LLMs achieve 22--66\% improvement over short-context baselines but lag human performance by 56\% on question answering and 73\% on temporal reasoning. RAG provides approximately 40\% improvement overall with particularly strong results on single-hop retrieval but struggles with temporal reasoning. Session decomposition improves retrieval by 4\% (Recall@k), and temporal-awareness yields 7--11\% improvement. These results underscore that conversational memory, particularly temporal and causal aspects, remains a significant unsolved challenge.

\subsubsection{MemoryBench and MemBench}

MemoryBench, accepted at ICLR~2025, provides the first benchmark offering all types of memory and feedback data for evaluating continual learning in LLM systems, addressing factual memory (information retention accuracy), reflective memory (understanding and summarization), and the role of user feedback in memory systems. MemBench evaluates memory from multiple perspectives---effectiveness (accuracy), efficiency (computational cost), and capacity (maximum storable information)---with findings that long-term memory shows 20--40\% accuracy degradation compared to short-term, computational cost increases up to 5$\times$ with memory size, and temporal reasoning improves 7--11\% with explicit temporal-awareness.

\subsubsection{MemoryAgentBench}

MemoryAgentBench evaluates long-term, interactive, and adaptive memory capabilities through incremental multi-turn interactions, assessing accurate retrieval, test-time learning (whether agents can learn during interaction), long-range understanding (relationships across many turns), and conflict resolution (handling contradictory information). These benchmarks collectively demonstrate that memory evaluation requires multi-dimensional assessment beyond simple accuracy, encompassing efficiency, capacity, temporal dynamics, and interaction-level adaptation.


% ============================================================
% SECTION 10: EVALUATION
% ============================================================
\section{Evaluation of Long-Context Capabilities}
\label{sec:evaluation}

Evaluating long-context capabilities requires testing not merely whether a model can accept long inputs but whether it effectively utilizes distributed information. The distinction between nominal context length (the maximum input length the model can process) and effective context length (the length at which the model maintains acceptable performance on meaningful tasks) is central to evaluation design. This section surveys the principal evaluation methodologies, from simple retrieval tests through comprehensive multi-task benchmarks, memory-specific metrics, and open problems in evaluation methodology.

\subsection{Needle-in-a-Haystack Testing}

The Needle-in-a-Haystack (NIAH) test \cite{kamradt2023needle} evaluates retrieval of specific information (the ``needle'') embedded at controlled positions within long irrelevant context (the ``haystack''), producing position-by-length accuracy heatmaps that visualize exactly where and when the model's retrieval capability breaks down. NIAH has become the standard screening test for long-context claims, with model developers routinely publishing NIAH heatmaps to demonstrate context window functionality.

However, NIAH has significant limitations that restrict its value as a standalone evaluation. The retrieval task is artificially simple---finding a single, distinctive piece of information in otherwise uniform text---bearing little resemblance to real-world information needs that require understanding relationships between multiple distributed pieces of information. Real tasks require multi-position integration rather than single-position retrieval. NIAH does not test generation quality, summarization capability, or reasoning over retrieved information. Multi-needle variants that require finding multiple pieces of information reveal 20--40\% accuracy drops compared to single-needle tests, and distractor variants that include plausible but incorrect information at other positions further reduce accuracy. These extended variants more closely approximate realistic use cases but remain synthetically structured. Despite these limitations, NIAH provides a necessary (though insufficient) condition for long-context capability: a model that fails NIAH cannot be expected to succeed on more demanding tasks.

\subsection{Comprehensive Benchmarks}

Several comprehensive benchmarks evaluate long-context capabilities across diverse task types, providing more realistic assessment than NIAH alone.

\subsubsection{RULER}

RULER \cite{hsieh2024ruler} goes beyond simple retrieval to evaluate four capability categories: multi-key retrieval (finding multiple pieces of information), variable tracking (following state changes across the context), aggregation (combining information from multiple positions), and multi-hop reasoning (performing inference chains across distributed information). RULER evaluates at controlled context lengths from 4K to 128K tokens, revealing that many models claiming 32K or larger context maintain acceptable performance on only about 50\% of challenging tasks. Multi-hop reasoning shows the steepest degradation, with most models performing near chance level at 128K tokens on complex reasoning tasks even when simple retrieval at the same length succeeds. This finding is particularly important because it demonstrates that the ability to find information does not imply the ability to reason about it.

\subsubsection{LongBench}

LongBench \cite{bai2024longbench} provides bilingual (English and Chinese), multitask evaluation across 21 datasets and six task categories: single-document QA, multi-document QA, summarization, few-shot learning, synthetic tasks, and code completion. Context lengths range from 6K to 20K tokens, placing LongBench in the medium-length evaluation range. The bilingual coverage is particularly valuable because most long-context benchmarks evaluate only English, leaving multilingual long-context capability largely uncharacterized.

\subsubsection{InfiniteBench}

InfiniteBench \cite{zhang2024infinitebench} extends evaluation beyond 100K tokens with tasks including book-length QA, novel character analysis, large codebase understanding, and mathematical reasoning over extended problem sets. Results reveal sobering performance at genuine long-context tasks: most models achieve only 20--50\% accuracy even on relatively simple tasks at these extreme lengths. InfiniteBench demonstrates that the gap between nominal and effective context length grows substantially at very long contexts, with no currently evaluated model approaching human-level performance at 200K+ tokens.

\subsubsection{BABILong}

BABILong \cite{kuratov2024babilong} embeds simple reasoning tasks from the original bAbI benchmark within contexts of up to 1~million tokens, providing controlled evaluation of reasoning at extreme lengths. The key finding is catastrophic: even simple reasoning tasks (counting, sorting, basic deduction) that models solve perfectly at short context lengths fail catastrophically beyond 32K tokens, with effective context utilization estimated at only 10--20\% of the nominal window. This result, combined with the finding that simple retrieval at the same lengths is often successful, indicates a fundamental gap between the ability to locate information and the ability to use it for reasoning---a gap that current architectures have not bridged.

\subsubsection{HELMET}

HELMET \cite{yen2024helmet} takes an application-centric approach, evaluating across seven task categories: retrieval, re-ranking, summarization, multi-document QA, many-shot in-context learning, long-document QA, and coding at context lengths from 4K to 128K tokens. A key finding is that model rankings vary significantly across task types and context lengths---a model that excels at retrieval may perform poorly at summarization, and a model that works well at 32K may degrade substantially at 128K. This result challenges the notion of a single ``long-context capability'' and suggests that evaluation must be multi-dimensional and task-specific.

\subsubsection{Additional Comprehensive Benchmarks}

L-Eval \cite{an2024leval} establishes standardized evaluation with closed-ended and open-ended tasks including code completion, providing automatic and human evaluation protocols. SCROLLS \cite{shaham2022scrolls} and ZeroSCROLLS \cite{shaham2023zeroscrolls} provide standardized evaluation over long sequences including legal documents and scientific papers in zero-shot settings. LongPPL proposes using perplexity-based metrics as efficient proxies for full task evaluation, with strong correlations ($r \approx 0.9$) between LongPPL scores and downstream task performance, enabling rapid model screening before expensive full evaluation.

\subsection{Domain-Specific and Specialized Benchmarks}

\subsubsection{Length-Stratified Evaluation}

LV-Eval provides balanced evaluation across five length levels (16K, 32K, 64K, 128K, and 256K tokens) with confusing facts insertion (adding plausible but incorrect information to test discrimination) and keyword-recall metrics that measure whether the model accesses the correct portion of context. Results reveal significant degradation at 128K+ tokens across all evaluated models, with confusing facts reducing accuracy by 10--25\% compared to clean contexts.

\subsubsection{Dependency-Aware Evaluation}

LooGLE evaluates long-dependency understanding with human-annotated questions specifically classified by the span of dependency required for the answer---how far apart in the text the relevant pieces of information are located. LooGLE finds 80--90\% accuracy on short dependencies (information within 1K tokens of each other) dropping to below 10\% for very long dependencies (information separated by 20K+ tokens), providing the most direct evidence of how context utilization degrades with information distance.

\subsubsection{RAG-Specific Evaluation}

The CRAG benchmark evaluates RAG systems across 4,409 QA pairs covering diverse domains and question types, revealing that state-of-the-art RAG solutions achieve approximately 63\% accuracy without hallucination---demonstrating both the promise and limitations of retrieval augmentation. Qasper evaluates long-document QA specifically on research papers, testing the ability to find and synthesize information within scientific documents. NarrativeQA tests comprehension over books and scripts averaging 60K tokens, requiring understanding of narrative structure, character development, and plot relationships.

\subsection{Memory-Specific Evaluation Metrics}

As LLMs transition from single-session to multi-session, long-horizon applications, memory-specific evaluation metrics are becoming increasingly important.

\subsubsection{Retrieval-Focused Metrics}

Standard information retrieval metrics apply to memory systems: Recall@k measures what fraction of relevant information appears in the top-$k$ retrieved items; Precision@k measures what fraction of retrieved items are relevant; Mean Reciprocal Rank (MRR) measures the average position of the first relevant result. Position-aware variants track how these metrics vary with memory age, revealing primacy and recency effects in memory retrieval systems.

\subsubsection{Temporal Metrics}

Temporal evaluation captures memory dynamics over time. Insertion position decay measures accuracy as a function of when information was stored in memory, quantifying how memory degrades with age. Typical findings show 95\% accuracy for immediately stored information dropping to 35\% for information stored 1,000+ steps ago. Temporal awareness measures improvement from explicit temporal information (timestamps, temporal references), with typical improvements of 7--11\%.

\subsubsection{Efficiency Metrics}

Memory efficiency metrics capture the computational cost of memory operations: tokens per operation (retrieval cost, update cost, summarization cost), latency of memory operations (critical for interactive applications), and the trade-off between memory capacity and retrieval precision. Memory systems show trade-offs between accuracy, efficiency, and capacity that must be evaluated jointly rather than optimizing any single dimension.

\subsection{Open Problems in Long-Context Evaluation}

\subsubsection{Scaling Laws for Long Context}

Scaling laws that predict performance based on compute, parameters, and training data have been foundational for LLM development, but long-context scaling laws remain poorly understood. Performance does not degrade smoothly with context length---cliff behaviors and phase transitions are observed at specific lengths. The optimal allocation of training data for long-context capability is unclear, with trade-offs between short-context and long-context performance that are not well characterized. Different architectures (Transformer, Mamba, hybrid) exhibit different degradation curves, and no unified framework predicts performance across architectures and lengths.

\subsubsection{Unified Evaluation Framework}

The proliferation of 10+ major benchmarks with inconsistent methodology, different metrics, different document sources, and different task definitions makes cross-benchmark comparison difficult. No single benchmark comprehensively evaluates all aspects of long-context capability, and results across benchmarks do not aggregate naturally. The evaluation field needs a consensus on minimal evaluation sets, standard metrics, and benchmark aggregation methods. The problem of benchmark gaming---models trained on test sets or similar data, leading to inflated performance claims---further complicates evaluation.

\subsubsection{Real-World Grounding}

Evaluation benchmarks are increasingly distant from actual deployment use cases. Synthetic tasks may not reflect real-world information needs. Real documents have characteristics (formatting, tables, code, multimodal content) not captured by clean benchmark texts. Query-document alignment differs between benchmarks (where queries are generated for documents) and real use (where documents are found for queries). Interactive, multi-turn evaluation (essential for conversational applications) is rarely tested in standard benchmarks.

\subsubsection{Efficient Evaluation Methodology}

Long-context evaluation is computationally expensive: NIAH takes minutes to hours per model, BABILong at 1M tokens takes days of computation, and InfiniteBench at 100K+ tokens requires weeks of evaluation for comprehensive assessment. This cost prohibits rapid iteration and limits evaluation to well-resourced organizations. Surrogate metrics like LongPPL offer promising efficiency gains but may miss important performance variations. The development of efficient, reliable evaluation protocols remains an active challenge.

\subsection{Evaluation Synthesis}

Comprehensive evaluation requires multiple complementary approaches: baseline screening (NIAH for basic retrieval capability), comprehensive benchmarks (LongBench, RULER, HELMET for multi-task assessment), domain-specific evaluation (Qasper for scientific documents, CRAG for RAG systems), memory evaluation (MemBench, LoCoMo for long-term memory), and efficiency assessment (latency, throughput, memory consumption). A critical meta-observation is the substantial gap between benchmark performance and real-world utilization: models achieving strong NIAH scores may fail on nuanced tasks, performance on controlled benchmarks may not predict deployment performance, and the evaluation methodology itself remains a subject of active research and debate.


% ============================================================
% SECTION 11: CONCLUSION
% ============================================================
\section{Conclusion and Future Directions}
\label{sec:conclusion}

This survey has examined context management in large language models through ten interconnected perspectives, from foundational attention mechanisms through human-inspired hierarchical memory architectures. The field has evolved rapidly from narrow technical solutions (extending context windows through positional encoding tricks) to broad architectural paradigms (hierarchical memory systems inspired by cognitive science). Several overarching themes and future directions emerge from this comprehensive review.

\subsection{The Convergence of Cognitive Science and Systems Design}

The most striking theme across the surveyed literature is the convergence between cognitive science principles and practical system design. This convergence manifests at multiple levels of the memory hierarchy. The Atkinson-Shiffrin model, with its three-store architecture of sensory, short-term, and long-term memory, provides an effective template for computational systems: the Hierarchical Memory Transformer implements a direct three-tier analogue, Strata applies the hierarchy to KV cache management in production, and MemGPT implements the hierarchy through explicit tiered storage with OS-inspired paging.

Baddeley's working memory model, with its central executive and specialized slave systems, informs architectures like CoALA where the language model serves as the central executive directing specialized memory operations. Tulving's distinction between episodic and semantic memory explains the asymmetry in current LLMs (strong semantic, weak episodic) and motivates explicit episodic memory systems like Generative Agents and EM-LLM. Ebbinghaus's forgetting curve informs memory management policies in MemoryBank and psychology-informed conversational agents. McClelland's complementary learning systems theory explains why purely parametric and purely retrieval-based approaches are each insufficient and why the most effective systems combine both, as demonstrated by RETRO, Self-RAG, and memory-augmented agent architectures.

This convergence is not coincidental: both biological and computational memory systems face the same fundamental trade-off between access speed and storage capacity, and hierarchical solutions address this through organizational principles equally applicable across substrates. The consistent success of cognitively-inspired architectures suggests that human memory organization reflects deep functional constraints on information processing systems rather than biological idiosyncrasies.

\subsection{From Extension to Architecture}

The field has progressed through three distinct phases. The first phase focused on extending existing architectures through position encoding innovations (PI, NTK, YaRN, LongRoPE) and efficient attention mechanisms (Longformer, BigBird, FlashAttention). These approaches increase the nominal context window while preserving the flat, unstructured token sequence as the fundamental representational unit.

The second phase introduced augmentation through external memory (RAG, MemGPT, Memorizing Transformer), adding retrieval mechanisms that supplement the context window without changing the model's architecture. These systems treat the context window as a fixed resource and manage it through intelligent loading and unloading of information.

The third and current phase involves fundamental rearchitecting of how models interact with information. Hierarchical memory systems (HMT, CoALA, Generative Agents) organize information across multiple temporal and abstraction scales. Self-directed memory management (MemGPT, SCM) gives the model agency over its own memory operations. Consolidation mechanisms (sleep-inspired replay, recursive summarization) enable knowledge to be transferred between memory systems over time. The most capable systems combine techniques from all three phases, using efficient attention within a context window managed by intelligent retrieval from hierarchical external memory---mirroring the multi-system biological memory organization that decades of cognitive science research has revealed.

\subsection{Open Challenges}

Despite the substantial progress surveyed in this review, several fundamental challenges remain unresolved.

Effective context utilization continues to lag far behind nominal context length. Models claiming 128K or million-token contexts show effective utilization of only 10--20\% of the nominal window on demanding tasks. The gap between finding information (NIAH-style retrieval) and using information (reasoning, synthesis, multi-hop inference) remains large and poorly understood. Closing this gap may require not just better architectures but better training procedures that expose models to genuine long-context reasoning tasks.

Temporal reasoning remains particularly difficult. Evaluations consistently show 50--73\% performance gaps between models and humans on temporal reasoning tasks. Conversational memory systems struggle to maintain accurate temporal representations across sessions, and the majority of current architectures lack explicit temporal reasoning mechanisms, relying on the language model's implicit understanding of temporal relationships.

Metacognitive calibration is essential for trustworthy memory management but remains poorly developed. Effective memory systems must know what they know (and do not know) to make good retrieval, storage, and compression decisions. Current models exhibit Judgments of Learning (JoL) that correlate only $r = 0.12$--$0.18$ with actual recall accuracy, compared to human calibration of $r \approx 0.68$. Without reliable metacognition, memory management decisions are poorly informed, leading to storing unimportant information, failing to retrieve relevant information, and over-confident generation without adequate evidential support.

Continual learning without catastrophic forgetting remains unsolved at scale. While parameter-efficient methods (LoRA) and rehearsal strategies (experience replay) mitigate forgetting in controlled settings, no current approach enables reliable continual knowledge acquisition over extended deployment lifetimes---a capability essential for persistent assistants that must integrate new information over months or years.

Unified evaluation methodology is needed. The proliferation of benchmarks with inconsistent methodology and metrics makes it difficult to assess genuine progress. The community lacks consensus on minimum evaluation standards, benchmark aggregation methods, and the relationship between benchmark performance and real-world deployment capability.

\subsection{Future Directions}

Several promising research directions emerge from the open challenges.

Sleep-inspired consolidation processes may provide principled mechanisms for maintaining memory quality over extended deployments. The success of NeuroDream (reducing forgetting from 40\% to 15\%) and sleep replay consolidation (preserving 86\% versus 18\% task accuracy) suggests that scheduled offline processing---reviewing, reorganizing, and consolidating accumulated memories---could substantially improve long-term memory quality. Production systems might implement ``maintenance windows'' during low-utilization periods, analogous to biological sleep, where accumulated episodic memories are consolidated into more permanent representations.

Importance-weighted encoding, inspired by emotional tagging in human memory, could enable systems to prioritize functionally important information during storage, compression, and retrieval. Current systems treat all tokens or memories with equal importance (or use simple heuristics like attention scores), but biological memory systems employ sophisticated importance evaluation that considers emotional salience, goal relevance, novelty, and encoding depth. Implementing analogous importance evaluation in LLM memory systems could dramatically improve the quality-efficiency trade-off in memory management.

Neuro-symbolic integration may provide the structured memory representations needed for complex multi-hop reasoning and temporal understanding. GraphRAG demonstrates that structured knowledge representations enable reasoning capabilities beyond what flat vector retrieval achieves. Extending this principle to full cognitive architectures that combine neural language understanding with symbolic knowledge representation and logical reasoning could address the reasoning gap that current benchmarks reveal.

Multimodal memory systems will become increasingly necessary as language models evolve toward general-purpose cognitive agents that process text, images, video, audio, and structured data. Current memory architectures are predominantly text-based, but real-world applications involve multimodal information that must be stored, retrieved, and reasoned about in an integrated manner.

Adaptive memory architectures that learn their own memory management policies---deciding what to store, when to retrieve, how to compress, and when to consolidate---represent perhaps the most promising long-term direction. Rather than hand-designing memory hierarchies and management rules, future systems may learn optimal memory strategies through reinforcement learning over extended interactions, discovering memory management principles that go beyond those identified by cognitive science.

\subsection{Closing Perspective}

The ultimate vision emerging from this survey is language model systems that manage information through a rich ecosystem of specialized, interacting memory systems---each optimized for its temporal scale and functional role, connected by consolidation processes that transfer knowledge between levels, and governed by metacognitive monitoring that assesses the reliability and relevance of stored information. The principles for building such systems are increasingly well-understood, drawn from decades of cognitive science research on human memory and implemented in the growing body of computational work surveyed here. The challenge ahead is integrating these principles into unified architectures that are simultaneously theoretically principled, practically deployable, and computationally scalable. The trajectory of the field---from simple context extension to cognitively-inspired hierarchical memory---suggests that the next generation of language model systems will be designed not merely as statistical pattern matchers with large input buffers, but as cognitive architectures with genuine memory management capabilities.

% ============================================================
% REFERENCES
% ============================================================
\bibliographystyle{IEEEtran}
\bibliography{references}

\end{document}
